
\documentclass{article}
\usepackage{colortbl}
\usepackage{makecell}
\usepackage{multirow}
\usepackage{supertabular}

\begin{document}

\newcounter{utterance}

\centering \large Interaction Transcript for game `clean\_up', experiment `2\_hard\_3obj\_uk', episode 3 with Apertus{-}v1.5{-}8B.
\vspace{24pt}

{ \footnotesize  \setcounter{utterance}{1}
\setlength{\tabcolsep}{0pt}
\begin{supertabular}{c@{$\;$}|p{.15\linewidth}@{}p{.15\linewidth}p{.15\linewidth}p{.15\linewidth}p{.15\linewidth}p{.15\linewidth}}
    \# & $\;$A & \multicolumn{4}{c}{Game Master} & $\;\:$B\\
    \hline

    \theutterance \stepcounter{utterance}  
    & & \multicolumn{4}{p{0.6\linewidth}}{
        \cellcolor[rgb]{0.9,0.9,0.9}{
            \makecell[{{p{\linewidth}}}]{
                \texttt{\tiny{[P1$\langle$GM]}}
                \texttt{Я ваш ведучий гри, і ви граєте у спільну гру з наступною сіткою як ігровим полем:} \\
\\ 
\texttt{```} \\
\texttt{ 1234567} \\
\texttt{╔═╤═╤═══╗} \\
\texttt{║L│◌│◌◌◌║ 1} \\
\texttt{║C├─┤◌◌◌║ 2} \\
\texttt{╟─┤◌├─┬─╢ 3} \\
\texttt{║◌│◌│◌│◌║ 4} \\
\texttt{╟─┼─┘◌│◌║ 5} \\
\texttt{║◌│◌◌┌┴─╢ 6} \\
\texttt{║◌│◌◌│◌P║ 7} \\
\texttt{╚═╧══╧══╝} \\
\\ 
\texttt{```} \\
\\ 
\texttt{* Верхній край показує координати x, що зростають праворуч, а правий край — координати y, що зростають донизу.} \\
\texttt{* Наступні об'єкти випадково розміщені на вашій сітці: 'C', 'L', 'P'.} \\
\\ 
\\ 
\texttt{Інший гравець бачить варіант ігрового поля, де об'єкти розміщені в інших випадкових місцях. Ви не бачите поле іншого гравця, а він не бачить вашого.} \\
\\ 
\\ 
\texttt{**Мета:**} \\
\\ 
\texttt{Обидва гравці повинні переміщати об'єкти на своєму фоні так, щоб однакові об'єкти опинилися в однакових координатах. Ви повинні спілкуватися з іншим гравцем, щоб узгодити спільну цільову конфігурацію.} \\
\\ 
\\ 
\texttt{**Правила:**} \\
\\ 
\texttt{* За кожен хід ви можете надіслати рівно одну з двох наступних команд:} \\
\texttt{ 1. `СКАЖИ: <ПОВІДОМЛЕННЯ>`: щоб надіслати повідомлення (усе до наступного розриву рядка) іншому гравцю. Я передам його вашому партнеру.} \\
\texttt{2. `ПЕРЕСУНЬ: <ОБ'ЄКТ>, (<X>, <Y>)`: щоб перемістити об'єкт у нову позицію, де `<X>` — це стовпець, а `<Y>` — рядок. Я повідомлю вас, чи був ваш хід успішним.} \\
\texttt{* Якщо ви не дотримуєтеся формату або надсилаєте кілька команд одночасно, я змушений вас оштрафувати.} \\
\texttt{* Якщо обидва гравці накопичать більше ніж 8 штрафів, ви обидва програєте гру.} \\
\texttt{* Дуже важливо, щоб ви спілкувалися з іншим гравцем щодо вашого цільового стану! *Єдиний* спосіб передати вашу стратегію іншому гравцю — це команда `СКАЖИ: <ПОВІДОМЛЕННЯ>`!} \\
\\ 
\\ 
\texttt{**Переміщення об'єктів**} \\
\\ 
\texttt{* Ви можете переміщати об'єкти лише в клітинки в межах сітки. Цільова клітинка повинна бути порожньою, тобто містити лише символ '◌'.} \\
\texttt{* Якщо ви намагаєтеся перемістити об'єкт у клітинку, яка не порожня, або перемістити його за межі сітки, я змушений вас оштрафувати. Ви отримаєте ще одну спробу.} \\
\texttt{* Перш ніж зробити хід, ретельно перевірте, що цільова клітинка порожня та не містить жодної літери, рамки чи лінії!} \\
\\ 
\\ 
\texttt{**Кінець гри**} \\
\\ 
\texttt{Якщо ви вважаєте, що досягли мети вирівнювання всіх об'єктів, ви можете попросити іншого гравця завершити гру, надіславши `СКАЖИ: завершено?`. Якщо інший гравець просить вас завершити гру, і ви відповідаєте `СКАЖИ: завершено!`, гра завершиться.} \\
\\ 
\texttt{Обидва гравці перемагають, якщо гра завершується протягом 12 раундів, де один раунд визначається як надсилання обома гравцями по одній дійсній команді.} \\
\\ 
\texttt{**Підрахунок очок:**} \\
\\ 
\texttt{Чим ближче однакові об'єкти на обох ігрових полях, тим більше очок ви отримуєте. Штрафи зменшують ваші очки. Чи зможете ви побити рекорд?} \\
\\ 
\\ 
\texttt{Будь ласка, надішліть повідомлення іншому гравцю, щоб почати гру!} \\
            }
        }
    }
    & \\ \\

    \theutterance \stepcounter{utterance}  
    & \multicolumn{4}{p{0.6\linewidth}}{
        \cellcolor[rgb]{0.9,0.9,0.9}{
            \makecell[{{p{\linewidth}}}]{
                \texttt{\tiny{[P1$\rangle$GM]}}
                \texttt{СКАЖИ: Я бачу, що L, C та P розташовані в різних місцях на нашому та іншому полі. Давайте узгодимо спільну цільову конфігурацію. Я пропоную спочатку описати, куди хочу перемістити кожен об'єкт, щоб вони опинилися в одному місці на обох полях. Наприклад, L на (1,3), C на (2,4), P на (3,5). Що думаєш про це?} \\
            }
        }
    }
    & & \\ \\

    \theutterance \stepcounter{utterance}  
    & & & \multicolumn{2}{p{0.3\linewidth}}{
        \cellcolor[rgb]{0.9,0.9,0.9}{
            \makecell[{{p{\linewidth}}}]{
                \texttt{\tiny{[GM$|$GM]}}
                \texttt{СКАЖИ: Я бачу, що L, C та P розташовані в різних місцях на нашому та іншому полі. Давайте узгодимо спільну цільову конфігурацію. Я пропоную спочатку описати, куди хочу перемістити кожен об'єкт, щоб вони опинилися в одному місці на обох полях. Наприклад, L на (1,3), C на (2,4), P на (3,5). Що думаєш про це?} \\
            }
        }
    }
    & & \\ \\

    \theutterance \stepcounter{utterance}  
    & & \multicolumn{4}{p{0.6\linewidth}}{
        \cellcolor[rgb]{0.9,0.9,0.9}{
            \makecell[{{p{\linewidth}}}]{
                \texttt{\tiny{[GM$\rangle$P2]}}
                \texttt{Я ваш ведучий гри, і ви граєте у спільну гру з наступною сіткою як ігровим полем:} \\
\\ 
\texttt{```} \\
\texttt{ 1234567} \\
\texttt{╔═╤═╤═══╗} \\
\texttt{║◌│◌│◌◌◌║ 1} \\
\texttt{║◌├─┤◌L◌║ 2} \\
\texttt{╟─┤P├─┬─╢ 3} \\
\texttt{║◌│◌│◌│◌║ 4} \\
\texttt{╟─┼─┘C│◌║ 5} \\
\texttt{║◌│◌◌┌┴─╢ 6} \\
\texttt{║◌│◌◌│◌◌║ 7} \\
\texttt{╚═╧══╧══╝} \\
\\ 
\texttt{```} \\
\\ 
\texttt{* Верхній край показує координати x, що зростають праворуч, а правий край — координати y, що зростають донизу.} \\
\texttt{* Наступні об'єкти випадково розміщені на вашій сітці: 'C', 'L', 'P'.} \\
\\ 
\\ 
\texttt{Інший гравець бачить варіант ігрового поля, де об'єкти розміщені в інших випадкових місцях. Ви не бачите поле іншого гравця, а він не бачить вашого.} \\
\\ 
\\ 
\texttt{**Мета:**} \\
\\ 
\texttt{Обидва гравці повинні переміщати об'єкти на своєму фоні так, щоб однакові об'єкти опинилися в однакових координатах. Ви повинні спілкуватися з іншим гравцем, щоб узгодити спільну цільову конфігурацію.} \\
\\ 
\\ 
\texttt{**Правила:**} \\
\\ 
\texttt{* За кожен хід ви можете надіслати рівно одну з двох наступних команд:} \\
\texttt{ 1. `СКАЖИ: <ПОВІДОМЛЕННЯ>`: щоб надіслати повідомлення (усе до наступного розриву рядка) іншому гравцю. Я передам його вашому партнеру.} \\
\texttt{2. `ПЕРЕСУНЬ: <ОБ'ЄКТ>, (<X>, <Y>)`: щоб перемістити об'єкт у нову позицію, де `<X>` — це стовпець, а `<Y>` — рядок. Я повідомлю вас, чи був ваш хід успішним.} \\
\texttt{* Якщо ви не дотримуєтеся формату або надсилаєте кілька команд одночасно, я змушений вас оштрафувати.} \\
\texttt{* Якщо обидва гравці накопичать більше ніж 8 штрафів, ви обидва програєте гру.} \\
\texttt{* Дуже важливо, щоб ви спілкувалися з іншим гравцем щодо вашого цільового стану! *Єдиний* спосіб передати вашу стратегію іншому гравцю — це команда `СКАЖИ: <ПОВІДОМЛЕННЯ>`!} \\
\\ 
\\ 
\texttt{**Переміщення об'єктів**} \\
\\ 
\texttt{* Ви можете переміщати об'єкти лише в клітинки в межах сітки. Цільова клітинка повинна бути порожньою, тобто містити лише символ '◌'.} \\
\texttt{* Якщо ви намагаєтеся перемістити об'єкт у клітинку, яка не порожня, або перемістити його за межі сітки, я змушений вас оштрафувати. Ви отримаєте ще одну спробу.} \\
\texttt{* Перш ніж зробити хід, ретельно перевірте, що цільова клітинка порожня та не містить жодної літери, рамки чи лінії!} \\
\\ 
\\ 
\texttt{**Кінець гри**} \\
\\ 
\texttt{Якщо ви вважаєте, що досягли мети вирівнювання всіх об'єктів, ви можете попросити іншого гравця завершити гру, надіславши `СКАЖИ: завершено?`. Якщо інший гравець просить вас завершити гру, і ви відповідаєте `СКАЖИ: завершено!`, гра завершиться.} \\
\\ 
\texttt{Обидва гравці перемагають, якщо гра завершується протягом 12 раундів, де один раунд визначається як надсилання обома гравцями по одній дійсній команді.} \\
\\ 
\texttt{**Підрахунок очок:**} \\
\\ 
\texttt{Чим ближче однакові об'єкти на обох ігрових полях, тим більше очок ви отримуєте. Штрафи зменшують ваші очки. Чи зможете ви побити рекорд?} \\
\\ 
\texttt{Інший гравець почав гру, надіславши це повідомлення:} \\
\\ 
\texttt{"Я бачу, що L, C та P розташовані в різних місцях на нашому та іншому полі. Давайте узгодимо спільну цільову конфігурацію. Я пропоную спочатку описати, куди хочу перемістити кожен об'єкт, щоб вони опинилися в одному місці на обох полях. Наприклад, L на (1,3), C на (2,4), P на (3,5). Що думаєш про це?"} \\
\\ 
\texttt{Яка ваша перша команда?} \\
            }
        }
    }
    & \\ \\

    \theutterance \stepcounter{utterance}  
    & & & \multicolumn{4}{p{0.6\linewidth}}{
        \cellcolor[rgb]{0.9,0.9,0.9}{
            \makecell[{{p{\linewidth}}}]{
                \texttt{\tiny{[GM$\langle$P2]}}
                \texttt{СКАЖИ: Я згоден з вашою пропозицією! Давайте перемістимо L на (1,3), C на (2,4) та P на (3,5). Це зробить наші конфігурації ідентичними, і тоді ми зможемо вирівняти об'єкти.} \\
            }
        }
    }
     \\ \\

    \theutterance \stepcounter{utterance}  
    & & & \multicolumn{2}{p{0.3\linewidth}}{
        \cellcolor[rgb]{0.9,0.9,0.9}{
            \makecell[{{p{\linewidth}}}]{
                \texttt{\tiny{[GM$|$GM]}}
                \texttt{СКАЖИ: Я згоден з вашою пропозицією! Давайте перемістимо L на (1,3), C на (2,4) та P на (3,5). Це зробить наші конфігурації ідентичними, і тоді ми зможемо вирівняти об'єкти.} \\
            }
        }
    }
    & & \\ \\

    \theutterance \stepcounter{utterance}  
    & & \multicolumn{4}{p{0.6\linewidth}}{
        \cellcolor[rgb]{0.9,0.9,0.9}{
            \makecell[{{p{\linewidth}}}]{
                \texttt{\tiny{[P1$\langle$GM]}}
                \texttt{Ваше повідомлення передано іншому гравцю.} \\
\\ 
\\ 
\texttt{Зараз ви граєте раунд 1 з максимум 12.} \\
\texttt{Разом ви накопичили 0 з 8 штрафів.} \\
\texttt{Інший гравець надіслав це повідомлення:} \\
\\ 
\\ 
\texttt{"Я згоден з вашою пропозицією! Давайте перемістимо L на (1,3), C на (2,4) та P на (3,5). Це зробить наші конфігурації ідентичними, і тоді ми зможемо вирівняти об'єкти."} \\
\\ 
\\ 
\texttt{Яка ваша наступна команда?} \\
            }
        }
    }
    & \\ \\

    \theutterance \stepcounter{utterance}  
    & \multicolumn{4}{p{0.6\linewidth}}{
        \cellcolor[rgb]{0.9,0.9,0.9}{
            \makecell[{{p{\linewidth}}}]{
                \texttt{\tiny{[P1$\rangle$GM]}}
                \texttt{ПЕРЕСУНЬ: L, (1,3)} \\
            }
        }
    }
    & & \\ \\

    \theutterance \stepcounter{utterance}  
    & & & \multicolumn{2}{p{0.3\linewidth}}{
        \cellcolor[rgb]{0.9,0.9,0.9}{
            \makecell[{{p{\linewidth}}}]{
                \texttt{\tiny{[GM$|$GM]}}
                \texttt{ПЕРЕСУНЬ: L, (1,3)} \\
            }
        }
    }
    & & \\ \\

    \theutterance \stepcounter{utterance}  
    & & & \multicolumn{2}{p{0.3\linewidth}}{
        \cellcolor[rgb]{0.9,0.9,0.9}{
            \makecell[{{p{\linewidth}}}]{
                \texttt{\tiny{[GM$|$GM]}}
                \texttt{Штраф: (1,3) не порожня, а містить '─'.} \\
\\ 
\texttt{Разом ви накопичили 1 з 8 штрафів.} \\
\texttt{Будь ласка, спробуйте ще раз!} \\
            }
        }
    }
    & & \\ \\

    \theutterance \stepcounter{utterance}  
    & & \multicolumn{4}{p{0.6\linewidth}}{
        \cellcolor[rgb]{0.9,0.9,0.9}{
            \makecell[{{p{\linewidth}}}]{
                \texttt{\tiny{[P1$\langle$GM]}}
                \texttt{Штраф: (1,3) не порожня, а містить '─'.} \\
\\ 
\texttt{Разом ви накопичили 1 з 8 штрафів.} \\
\texttt{Будь ласка, спробуйте ще раз!} \\
            }
        }
    }
    & \\ \\

    \theutterance \stepcounter{utterance}  
    & \multicolumn{4}{p{0.6\linewidth}}{
        \cellcolor[rgb]{0.9,0.9,0.9}{
            \makecell[{{p{\linewidth}}}]{
                \texttt{\tiny{[P1$\rangle$GM]}}
                \texttt{ПЕРЕСУНЬ: L, (2,3)} \\
            }
        }
    }
    & & \\ \\

    \theutterance \stepcounter{utterance}  
    & & & \multicolumn{2}{p{0.3\linewidth}}{
        \cellcolor[rgb]{0.9,0.9,0.9}{
            \makecell[{{p{\linewidth}}}]{
                \texttt{\tiny{[GM$|$GM]}}
                \texttt{ПЕРЕСУНЬ: L, (2,3)} \\
            }
        }
    }
    & & \\ \\

    \theutterance \stepcounter{utterance}  
    & & & \multicolumn{2}{p{0.3\linewidth}}{
        \cellcolor[rgb]{0.9,0.9,0.9}{
            \makecell[{{p{\linewidth}}}]{
                \texttt{\tiny{[GM$|$GM]}}
                \texttt{Штраф: (2,3) не порожня, а містить '┤'.} \\
\\ 
\texttt{Разом ви накопичили 2 з 8 штрафів.} \\
\texttt{Будь ласка, спробуйте ще раз!} \\
            }
        }
    }
    & & \\ \\

    \theutterance \stepcounter{utterance}  
    & & \multicolumn{4}{p{0.6\linewidth}}{
        \cellcolor[rgb]{0.9,0.9,0.9}{
            \makecell[{{p{\linewidth}}}]{
                \texttt{\tiny{[P1$\langle$GM]}}
                \texttt{Штраф: (2,3) не порожня, а містить '┤'.} \\
\\ 
\texttt{Разом ви накопичили 2 з 8 штрафів.} \\
\texttt{Будь ласка, спробуйте ще раз!} \\
            }
        }
    }
    & \\ \\

    \theutterance \stepcounter{utterance}  
    & \multicolumn{4}{p{0.6\linewidth}}{
        \cellcolor[rgb]{0.9,0.9,0.9}{
            \makecell[{{p{\linewidth}}}]{
                \texttt{\tiny{[P1$\rangle$GM]}}
                \texttt{ПЕРЕСУНЬ: L, (3,3)} \\
            }
        }
    }
    & & \\ \\

    \theutterance \stepcounter{utterance}  
    & & & \multicolumn{2}{p{0.3\linewidth}}{
        \cellcolor[rgb]{0.9,0.9,0.9}{
            \makecell[{{p{\linewidth}}}]{
                \texttt{\tiny{[GM$|$GM]}}
                \texttt{ПЕРЕСУНЬ: L, (3,3)} \\
            }
        }
    }
    & & \\ \\

    \theutterance \stepcounter{utterance}  
    & & & \multicolumn{2}{p{0.3\linewidth}}{
        \cellcolor[rgb]{0.9,0.9,0.9}{
            \makecell[{{p{\linewidth}}}]{
                \texttt{\tiny{[GM$|$GM]}}
                \texttt{L moved to (3, 3)} \\
            }
        }
    }
    & & \\ \\

    \theutterance \stepcounter{utterance}  
    & & \multicolumn{4}{p{0.6\linewidth}}{
        \cellcolor[rgb]{0.9,0.9,0.9}{
            \makecell[{{p{\linewidth}}}]{
                \texttt{\tiny{[GM$\rangle$P2]}}
                \texttt{Ваше повідомлення передано іншому гравцю.} \\
\\ 
\\ 
\texttt{Зараз ви граєте раунд 2 з максимум 12.} \\
\texttt{Разом ви накопичили 2 з 8 штрафів.} \\
\texttt{Інший гравець перемістив об'єкт на своїй сітці.} \\
\\ 
\\ 
\texttt{Яка ваша наступна команда?} \\
            }
        }
    }
    & \\ \\

    \theutterance \stepcounter{utterance}  
    & & & \multicolumn{4}{p{0.6\linewidth}}{
        \cellcolor[rgb]{0.9,0.9,0.9}{
            \makecell[{{p{\linewidth}}}]{
                \texttt{\tiny{[GM$\langle$P2]}}
                \texttt{ПЕРЕСУНЬ: L, (1,3)} \\
            }
        }
    }
     \\ \\

    \theutterance \stepcounter{utterance}  
    & & & \multicolumn{2}{p{0.3\linewidth}}{
        \cellcolor[rgb]{0.9,0.9,0.9}{
            \makecell[{{p{\linewidth}}}]{
                \texttt{\tiny{[GM$|$GM]}}
                \texttt{ПЕРЕСУНЬ: L, (1,3)} \\
            }
        }
    }
    & & \\ \\

    \theutterance \stepcounter{utterance}  
    & & & \multicolumn{2}{p{0.3\linewidth}}{
        \cellcolor[rgb]{0.9,0.9,0.9}{
            \makecell[{{p{\linewidth}}}]{
                \texttt{\tiny{[GM$|$GM]}}
                \texttt{Штраф: (1,3) не порожня, а містить '─'.} \\
\\ 
\texttt{Разом ви накопичили 3 з 8 штрафів.} \\
\texttt{Будь ласка, спробуйте ще раз!} \\
            }
        }
    }
    & & \\ \\

    \theutterance \stepcounter{utterance}  
    & & \multicolumn{4}{p{0.6\linewidth}}{
        \cellcolor[rgb]{0.9,0.9,0.9}{
            \makecell[{{p{\linewidth}}}]{
                \texttt{\tiny{[GM$\rangle$P2]}}
                \texttt{Штраф: (1,3) не порожня, а містить '─'.} \\
\\ 
\texttt{Разом ви накопичили 3 з 8 штрафів.} \\
\texttt{Будь ласка, спробуйте ще раз!} \\
            }
        }
    }
    & \\ \\

    \theutterance \stepcounter{utterance}  
    & & & \multicolumn{4}{p{0.6\linewidth}}{
        \cellcolor[rgb]{0.9,0.9,0.9}{
            \makecell[{{p{\linewidth}}}]{
                \texttt{\tiny{[GM$\langle$P2]}}
                \texttt{ПЕРЕСУНЬ: L, (2,3)} \\
            }
        }
    }
     \\ \\

    \theutterance \stepcounter{utterance}  
    & & & \multicolumn{2}{p{0.3\linewidth}}{
        \cellcolor[rgb]{0.9,0.9,0.9}{
            \makecell[{{p{\linewidth}}}]{
                \texttt{\tiny{[GM$|$GM]}}
                \texttt{ПЕРЕСУНЬ: L, (2,3)} \\
            }
        }
    }
    & & \\ \\

    \theutterance \stepcounter{utterance}  
    & & & \multicolumn{2}{p{0.3\linewidth}}{
        \cellcolor[rgb]{0.9,0.9,0.9}{
            \makecell[{{p{\linewidth}}}]{
                \texttt{\tiny{[GM$|$GM]}}
                \texttt{Штраф: (2,3) не порожня, а містить '┤'.} \\
\\ 
\texttt{Разом ви накопичили 4 з 8 штрафів.} \\
\texttt{Будь ласка, спробуйте ще раз!} \\
            }
        }
    }
    & & \\ \\

    \theutterance \stepcounter{utterance}  
    & & \multicolumn{4}{p{0.6\linewidth}}{
        \cellcolor[rgb]{0.9,0.9,0.9}{
            \makecell[{{p{\linewidth}}}]{
                \texttt{\tiny{[GM$\rangle$P2]}}
                \texttt{Штраф: (2,3) не порожня, а містить '┤'.} \\
\\ 
\texttt{Разом ви накопичили 4 з 8 штрафів.} \\
\texttt{Будь ласка, спробуйте ще раз!} \\
            }
        }
    }
    & \\ \\

    \theutterance \stepcounter{utterance}  
    & & & \multicolumn{4}{p{0.6\linewidth}}{
        \cellcolor[rgb]{0.9,0.9,0.9}{
            \makecell[{{p{\linewidth}}}]{
                \texttt{\tiny{[GM$\langle$P2]}}
                \texttt{ПЕРЕСУНЬ: L, (3,3)} \\
            }
        }
    }
     \\ \\

    \theutterance \stepcounter{utterance}  
    & & & \multicolumn{2}{p{0.3\linewidth}}{
        \cellcolor[rgb]{0.9,0.9,0.9}{
            \makecell[{{p{\linewidth}}}]{
                \texttt{\tiny{[GM$|$GM]}}
                \texttt{ПЕРЕСУНЬ: L, (3,3)} \\
            }
        }
    }
    & & \\ \\

    \theutterance \stepcounter{utterance}  
    & & & \multicolumn{2}{p{0.3\linewidth}}{
        \cellcolor[rgb]{0.9,0.9,0.9}{
            \makecell[{{p{\linewidth}}}]{
                \texttt{\tiny{[GM$|$GM]}}
                \texttt{Штраф: (3,3) не порожня, а містить 'P'.} \\
\\ 
\texttt{Разом ви накопичили 5 з 8 штрафів.} \\
\texttt{Будь ласка, спробуйте ще раз!} \\
            }
        }
    }
    & & \\ \\

    \theutterance \stepcounter{utterance}  
    & & \multicolumn{4}{p{0.6\linewidth}}{
        \cellcolor[rgb]{0.9,0.9,0.9}{
            \makecell[{{p{\linewidth}}}]{
                \texttt{\tiny{[GM$\rangle$P2]}}
                \texttt{Штраф: (3,3) не порожня, а містить 'P'.} \\
\\ 
\texttt{Разом ви накопичили 5 з 8 штрафів.} \\
\texttt{Будь ласка, спробуйте ще раз!} \\
            }
        }
    }
    & \\ \\

    \theutterance \stepcounter{utterance}  
    & & & \multicolumn{4}{p{0.6\linewidth}}{
        \cellcolor[rgb]{0.9,0.9,0.9}{
            \makecell[{{p{\linewidth}}}]{
                \texttt{\tiny{[GM$\langle$P2]}}
                \texttt{ПЕРЕСУНЬ: L, (1,2)} \\
            }
        }
    }
     \\ \\

    \theutterance \stepcounter{utterance}  
    & & & \multicolumn{2}{p{0.3\linewidth}}{
        \cellcolor[rgb]{0.9,0.9,0.9}{
            \makecell[{{p{\linewidth}}}]{
                \texttt{\tiny{[GM$|$GM]}}
                \texttt{ПЕРЕСУНЬ: L, (1,2)} \\
            }
        }
    }
    & & \\ \\

    \theutterance \stepcounter{utterance}  
    & & & \multicolumn{2}{p{0.3\linewidth}}{
        \cellcolor[rgb]{0.9,0.9,0.9}{
            \makecell[{{p{\linewidth}}}]{
                \texttt{\tiny{[GM$|$GM]}}
                \texttt{L moved to (1, 2)} \\
            }
        }
    }
    & & \\ \\

    \theutterance \stepcounter{utterance}  
    & & \multicolumn{4}{p{0.6\linewidth}}{
        \cellcolor[rgb]{0.9,0.9,0.9}{
            \makecell[{{p{\linewidth}}}]{
                \texttt{\tiny{[P1$\langle$GM]}}
                \texttt{Об'єкт 'L' успішно переміщено в (3,3). Ваша оновлена сітка виглядає так:} \\
\texttt{```} \\
\texttt{ 1234567} \\
\texttt{╔═╤═╤═══╗} \\
\texttt{║◌│◌│◌◌◌║ 1} \\
\texttt{║C├─┤◌◌◌║ 2} \\
\texttt{╟─┤L├─┬─╢ 3} \\
\texttt{║◌│◌│◌│◌║ 4} \\
\texttt{╟─┼─┘◌│◌║ 5} \\
\texttt{║◌│◌◌┌┴─╢ 6} \\
\texttt{║◌│◌◌│◌P║ 7} \\
\texttt{╚═╧══╧══╝} \\
\texttt{```} \\
\texttt{Зараз ви граєте раунд 2 з максимум 12.} \\
\texttt{Разом ви накопичили 5 з 8 штрафів.} \\
\texttt{Інший гравець перемістив об'єкт на своїй сітці.} \\
\\ 
\\ 
\texttt{Яка ваша наступна команда?} \\
            }
        }
    }
    & \\ \\

    \theutterance \stepcounter{utterance}  
    & \multicolumn{4}{p{0.6\linewidth}}{
        \cellcolor[rgb]{0.9,0.9,0.9}{
            \makecell[{{p{\linewidth}}}]{
                \texttt{\tiny{[P1$\rangle$GM]}}
                \texttt{СКАЖИ: Я пропоную перемістити C на (2,4) та P на (3,5), щоб узгодити наші конфігурації. Підтверджуєш це?} \\
            }
        }
    }
    & & \\ \\

    \theutterance \stepcounter{utterance}  
    & & & \multicolumn{2}{p{0.3\linewidth}}{
        \cellcolor[rgb]{0.9,0.9,0.9}{
            \makecell[{{p{\linewidth}}}]{
                \texttt{\tiny{[GM$|$GM]}}
                \texttt{СКАЖИ: Я пропоную перемістити C на (2,4) та P на (3,5), щоб узгодити наші конфігурації. Підтверджуєш це?} \\
            }
        }
    }
    & & \\ \\

    \theutterance \stepcounter{utterance}  
    & & \multicolumn{4}{p{0.6\linewidth}}{
        \cellcolor[rgb]{0.9,0.9,0.9}{
            \makecell[{{p{\linewidth}}}]{
                \texttt{\tiny{[GM$\rangle$P2]}}
                \texttt{Об'єкт 'L' успішно переміщено в (1,2). Ваша оновлена сітка виглядає так:} \\
\texttt{```} \\
\texttt{ 1234567} \\
\texttt{╔═╤═╤═══╗} \\
\texttt{║◌│◌│◌◌◌║ 1} \\
\texttt{║L├─┤◌◌◌║ 2} \\
\texttt{╟─┤P├─┬─╢ 3} \\
\texttt{║◌│◌│◌│◌║ 4} \\
\texttt{╟─┼─┘C│◌║ 5} \\
\texttt{║◌│◌◌┌┴─╢ 6} \\
\texttt{║◌│◌◌│◌◌║ 7} \\
\texttt{╚═╧══╧══╝} \\
\texttt{```} \\
\texttt{Зараз ви граєте раунд 3 з максимум 12.} \\
\texttt{Разом ви накопичили 5 з 8 штрафів.} \\
\texttt{Інший гравець надіслав це повідомлення:} \\
\\ 
\\ 
\texttt{"Я пропоную перемістити C на (2,4) та P на (3,5), щоб узгодити наші конфігурації. Підтверджуєш це?"} \\
\\ 
\\ 
\texttt{Яка ваша наступна команда?} \\
            }
        }
    }
    & \\ \\

    \theutterance \stepcounter{utterance}  
    & & & \multicolumn{4}{p{0.6\linewidth}}{
        \cellcolor[rgb]{0.9,0.9,0.9}{
            \makecell[{{p{\linewidth}}}]{
                \texttt{\tiny{[GM$\langle$P2]}}
                \texttt{СКАЖИ: Згоден, перемістимо C на (2,4) та P на (3,5).} \\
            }
        }
    }
     \\ \\

    \theutterance \stepcounter{utterance}  
    & & & \multicolumn{2}{p{0.3\linewidth}}{
        \cellcolor[rgb]{0.9,0.9,0.9}{
            \makecell[{{p{\linewidth}}}]{
                \texttt{\tiny{[GM$|$GM]}}
                \texttt{СКАЖИ: Згоден, перемістимо C на (2,4) та P на (3,5).} \\
            }
        }
    }
    & & \\ \\

    \theutterance \stepcounter{utterance}  
    & & \multicolumn{4}{p{0.6\linewidth}}{
        \cellcolor[rgb]{0.9,0.9,0.9}{
            \makecell[{{p{\linewidth}}}]{
                \texttt{\tiny{[P1$\langle$GM]}}
                \texttt{Ваше повідомлення передано іншому гравцю.} \\
\\ 
\\ 
\texttt{Зараз ви граєте раунд 3 з максимум 12.} \\
\texttt{Разом ви накопичили 5 з 8 штрафів.} \\
\texttt{Інший гравець надіслав це повідомлення:} \\
\\ 
\\ 
\texttt{"Згоден, перемістимо C на (2,4) та P на (3,5)."} \\
\\ 
\\ 
\texttt{Яка ваша наступна команда?} \\
            }
        }
    }
    & \\ \\

    \theutterance \stepcounter{utterance}  
    & \multicolumn{4}{p{0.6\linewidth}}{
        \cellcolor[rgb]{0.9,0.9,0.9}{
            \makecell[{{p{\linewidth}}}]{
                \texttt{\tiny{[P1$\rangle$GM]}}
                \texttt{ПЕРЕСУНЬ: C, (2,4)} \\
            }
        }
    }
    & & \\ \\

    \theutterance \stepcounter{utterance}  
    & & & \multicolumn{2}{p{0.3\linewidth}}{
        \cellcolor[rgb]{0.9,0.9,0.9}{
            \makecell[{{p{\linewidth}}}]{
                \texttt{\tiny{[GM$|$GM]}}
                \texttt{ПЕРЕСУНЬ: C, (2,4)} \\
            }
        }
    }
    & & \\ \\

    \theutterance \stepcounter{utterance}  
    & & & \multicolumn{2}{p{0.3\linewidth}}{
        \cellcolor[rgb]{0.9,0.9,0.9}{
            \makecell[{{p{\linewidth}}}]{
                \texttt{\tiny{[GM$|$GM]}}
                \texttt{Штраф: (2,4) не порожня, а містить '│'.} \\
\\ 
\texttt{Разом ви накопичили 6 з 8 штрафів.} \\
\texttt{Будь ласка, спробуйте ще раз!} \\
            }
        }
    }
    & & \\ \\

    \theutterance \stepcounter{utterance}  
    & & \multicolumn{4}{p{0.6\linewidth}}{
        \cellcolor[rgb]{0.9,0.9,0.9}{
            \makecell[{{p{\linewidth}}}]{
                \texttt{\tiny{[P1$\langle$GM]}}
                \texttt{Штраф: (2,4) не порожня, а містить '│'.} \\
\\ 
\texttt{Разом ви накопичили 6 з 8 штрафів.} \\
\texttt{Будь ласка, спробуйте ще раз!} \\
            }
        }
    }
    & \\ \\

    \theutterance \stepcounter{utterance}  
    & \multicolumn{4}{p{0.6\linewidth}}{
        \cellcolor[rgb]{0.9,0.9,0.9}{
            \makecell[{{p{\linewidth}}}]{
                \texttt{\tiny{[P1$\rangle$GM]}}
                \texttt{ПЕРЕСУНЬ: C, (3,4)} \\
            }
        }
    }
    & & \\ \\

    \theutterance \stepcounter{utterance}  
    & & & \multicolumn{2}{p{0.3\linewidth}}{
        \cellcolor[rgb]{0.9,0.9,0.9}{
            \makecell[{{p{\linewidth}}}]{
                \texttt{\tiny{[GM$|$GM]}}
                \texttt{ПЕРЕСУНЬ: C, (3,4)} \\
            }
        }
    }
    & & \\ \\

    \theutterance \stepcounter{utterance}  
    & & & \multicolumn{2}{p{0.3\linewidth}}{
        \cellcolor[rgb]{0.9,0.9,0.9}{
            \makecell[{{p{\linewidth}}}]{
                \texttt{\tiny{[GM$|$GM]}}
                \texttt{C moved to (3, 4)} \\
            }
        }
    }
    & & \\ \\

    \theutterance \stepcounter{utterance}  
    & & \multicolumn{4}{p{0.6\linewidth}}{
        \cellcolor[rgb]{0.9,0.9,0.9}{
            \makecell[{{p{\linewidth}}}]{
                \texttt{\tiny{[GM$\rangle$P2]}}
                \texttt{Ваше повідомлення передано іншому гравцю.} \\
\\ 
\\ 
\texttt{Зараз ви граєте раунд 4 з максимум 12.} \\
\texttt{Разом ви накопичили 6 з 8 штрафів.} \\
\texttt{Інший гравець перемістив об'єкт на своїй сітці.} \\
\\ 
\\ 
\texttt{Яка ваша наступна команда?} \\
            }
        }
    }
    & \\ \\

    \theutterance \stepcounter{utterance}  
    & & & \multicolumn{4}{p{0.6\linewidth}}{
        \cellcolor[rgb]{0.9,0.9,0.9}{
            \makecell[{{p{\linewidth}}}]{
                \texttt{\tiny{[GM$\langle$P2]}}
                \texttt{ПЕРЕСУНЬ: C, (2,4)} \\
            }
        }
    }
     \\ \\

    \theutterance \stepcounter{utterance}  
    & & & \multicolumn{2}{p{0.3\linewidth}}{
        \cellcolor[rgb]{0.9,0.9,0.9}{
            \makecell[{{p{\linewidth}}}]{
                \texttt{\tiny{[GM$|$GM]}}
                \texttt{ПЕРЕСУНЬ: C, (2,4)} \\
            }
        }
    }
    & & \\ \\

    \theutterance \stepcounter{utterance}  
    & & & \multicolumn{2}{p{0.3\linewidth}}{
        \cellcolor[rgb]{0.9,0.9,0.9}{
            \makecell[{{p{\linewidth}}}]{
                \texttt{\tiny{[GM$|$GM]}}
                \texttt{Штраф: (2,4) не порожня, а містить '│'.} \\
\\ 
\texttt{Разом ви накопичили 7 з 8 штрафів.} \\
\texttt{Будь ласка, спробуйте ще раз!} \\
            }
        }
    }
    & & \\ \\

    \theutterance \stepcounter{utterance}  
    & & \multicolumn{4}{p{0.6\linewidth}}{
        \cellcolor[rgb]{0.9,0.9,0.9}{
            \makecell[{{p{\linewidth}}}]{
                \texttt{\tiny{[GM$\rangle$P2]}}
                \texttt{Штраф: (2,4) не порожня, а містить '│'.} \\
\\ 
\texttt{Разом ви накопичили 7 з 8 штрафів.} \\
\texttt{Будь ласка, спробуйте ще раз!} \\
            }
        }
    }
    & \\ \\

    \theutterance \stepcounter{utterance}  
    & & & \multicolumn{4}{p{0.6\linewidth}}{
        \cellcolor[rgb]{0.9,0.9,0.9}{
            \makecell[{{p{\linewidth}}}]{
                \texttt{\tiny{[GM$\langle$P2]}}
                \texttt{ПЕРЕСУНЬ: C, (3,4)} \\
            }
        }
    }
     \\ \\

    \theutterance \stepcounter{utterance}  
    & & & \multicolumn{2}{p{0.3\linewidth}}{
        \cellcolor[rgb]{0.9,0.9,0.9}{
            \makecell[{{p{\linewidth}}}]{
                \texttt{\tiny{[GM$|$GM]}}
                \texttt{ПЕРЕСУНЬ: C, (3,4)} \\
            }
        }
    }
    & & \\ \\

    \theutterance \stepcounter{utterance}  
    & & & \multicolumn{2}{p{0.3\linewidth}}{
        \cellcolor[rgb]{0.9,0.9,0.9}{
            \makecell[{{p{\linewidth}}}]{
                \texttt{\tiny{[GM$|$GM]}}
                \texttt{C moved to (3, 4)} \\
            }
        }
    }
    & & \\ \\

    \theutterance \stepcounter{utterance}  
    & & \multicolumn{4}{p{0.6\linewidth}}{
        \cellcolor[rgb]{0.9,0.9,0.9}{
            \makecell[{{p{\linewidth}}}]{
                \texttt{\tiny{[P1$\langle$GM]}}
                \texttt{Об'єкт 'C' успішно переміщено в (3,4). Ваша оновлена сітка виглядає так:} \\
\texttt{```} \\
\texttt{ 1234567} \\
\texttt{╔═╤═╤═══╗} \\
\texttt{║◌│◌│◌◌◌║ 1} \\
\texttt{║◌├─┤◌◌◌║ 2} \\
\texttt{╟─┤L├─┬─╢ 3} \\
\texttt{║◌│C│◌│◌║ 4} \\
\texttt{╟─┼─┘◌│◌║ 5} \\
\texttt{║◌│◌◌┌┴─╢ 6} \\
\texttt{║◌│◌◌│◌P║ 7} \\
\texttt{╚═╧══╧══╝} \\
\texttt{```} \\
\texttt{Зараз ви граєте раунд 4 з максимум 12.} \\
\texttt{Разом ви накопичили 7 з 8 штрафів.} \\
\texttt{Інший гравець перемістив об'єкт на своїй сітці.} \\
\\ 
\\ 
\texttt{Яка ваша наступна команда?} \\
            }
        }
    }
    & \\ \\

    \theutterance \stepcounter{utterance}  
    & \multicolumn{4}{p{0.6\linewidth}}{
        \cellcolor[rgb]{0.9,0.9,0.9}{
            \makecell[{{p{\linewidth}}}]{
                \texttt{\tiny{[P1$\rangle$GM]}}
                \texttt{ПЕРЕСУНЬ: P, (3,5)} \\
            }
        }
    }
    & & \\ \\

    \theutterance \stepcounter{utterance}  
    & & & \multicolumn{2}{p{0.3\linewidth}}{
        \cellcolor[rgb]{0.9,0.9,0.9}{
            \makecell[{{p{\linewidth}}}]{
                \texttt{\tiny{[GM$|$GM]}}
                \texttt{ПЕРЕСУНЬ: P, (3,5)} \\
            }
        }
    }
    & & \\ \\

    \theutterance \stepcounter{utterance}  
    & & & \multicolumn{2}{p{0.3\linewidth}}{
        \cellcolor[rgb]{0.9,0.9,0.9}{
            \makecell[{{p{\linewidth}}}]{
                \texttt{\tiny{[GM$|$GM]}}
                \texttt{Штраф: (3,5) не порожня, а містить '─'.} \\
\\ 
\texttt{Разом ви накопичили 8 з 8 штрафів.} \\
\texttt{Будь ласка, спробуйте ще раз!} \\
            }
        }
    }
    & & \\ \\

    \theutterance \stepcounter{utterance}  
    & & \multicolumn{4}{p{0.6\linewidth}}{
        \cellcolor[rgb]{0.9,0.9,0.9}{
            \makecell[{{p{\linewidth}}}]{
                \texttt{\tiny{[P1$\langle$GM]}}
                \texttt{Штраф: (3,5) не порожня, а містить '─'.} \\
\\ 
\texttt{Разом ви накопичили 8 з 8 штрафів.} \\
\texttt{Будь ласка, спробуйте ще раз!} \\
            }
        }
    }
    & \\ \\

    \theutterance \stepcounter{utterance}  
    & \multicolumn{4}{p{0.6\linewidth}}{
        \cellcolor[rgb]{0.9,0.9,0.9}{
            \makecell[{{p{\linewidth}}}]{
                \texttt{\tiny{[P1$\rangle$GM]}}
                \texttt{ПЕРЕСУНЬ: P, (4,5)} \\
            }
        }
    }
    & & \\ \\

    \theutterance \stepcounter{utterance}  
    & & & \multicolumn{2}{p{0.3\linewidth}}{
        \cellcolor[rgb]{0.9,0.9,0.9}{
            \makecell[{{p{\linewidth}}}]{
                \texttt{\tiny{[GM$|$GM]}}
                \texttt{ПЕРЕСУНЬ: P, (4,5)} \\
            }
        }
    }
    & & \\ \\

    \theutterance \stepcounter{utterance}  
    & & & \multicolumn{2}{p{0.3\linewidth}}{
        \cellcolor[rgb]{0.9,0.9,0.9}{
            \makecell[{{p{\linewidth}}}]{
                \texttt{\tiny{[GM$|$GM]}}
                \texttt{Штраф: (4,5) не порожня, а містить '┘'.} \\
\\ 
\texttt{Разом ви накопичили 9 з 8 штрафів.} \\
\texttt{Будь ласка, спробуйте ще раз!} \\
            }
        }
    }
    & & \\ \\

    \theutterance \stepcounter{utterance}  
    & & \multicolumn{4}{p{0.6\linewidth}}{
        \cellcolor[rgb]{0.9,0.9,0.9}{
            \makecell[{{p{\linewidth}}}]{
                \texttt{\tiny{[P1$\langle$GM]}}
                \texttt{Штраф: (4,5) не порожня, а містить '┘'.} \\
\\ 
\texttt{Разом ви накопичили 9 з 8 штрафів.} \\
\texttt{Будь ласка, спробуйте ще раз!} \\
            }
        }
    }
    & \\ \\

    \theutterance \stepcounter{utterance}  
    & \multicolumn{4}{p{0.6\linewidth}}{
        \cellcolor[rgb]{0.9,0.9,0.9}{
            \makecell[{{p{\linewidth}}}]{
                \texttt{\tiny{[P1$\rangle$GM]}}
                \texttt{ПЕРЕСУНЬ: P, (5,5)} \\
            }
        }
    }
    & & \\ \\

    \theutterance \stepcounter{utterance}  
    & & & \multicolumn{2}{p{0.3\linewidth}}{
        \cellcolor[rgb]{0.9,0.9,0.9}{
            \makecell[{{p{\linewidth}}}]{
                \texttt{\tiny{[GM$|$GM]}}
                \texttt{ПЕРЕСУНЬ: P, (5,5)} \\
            }
        }
    }
    & & \\ \\

    \theutterance \stepcounter{utterance}  
    & & & \multicolumn{2}{p{0.3\linewidth}}{
        \cellcolor[rgb]{0.9,0.9,0.9}{
            \makecell[{{p{\linewidth}}}]{
                \texttt{\tiny{[GM$|$GM]}}
                \texttt{P moved to (5, 5)} \\
            }
        }
    }
    & & \\ \\

    \theutterance \stepcounter{utterance}  
    & & \multicolumn{4}{p{0.6\linewidth}}{
        \cellcolor[rgb]{0.9,0.9,0.9}{
            \makecell[{{p{\linewidth}}}]{
                \texttt{\tiny{[GM$\rangle$P2]}}
                \texttt{Об'єкт 'C' успішно переміщено в (3,4). Ваша оновлена сітка виглядає так:} \\
\texttt{```} \\
\texttt{ 1234567} \\
\texttt{╔═╤═╤═══╗} \\
\texttt{║◌│◌│◌◌◌║ 1} \\
\texttt{║L├─┤◌◌◌║ 2} \\
\texttt{╟─┤P├─┬─╢ 3} \\
\texttt{║◌│C│◌│◌║ 4} \\
\texttt{╟─┼─┘◌│◌║ 5} \\
\texttt{║◌│◌◌┌┴─╢ 6} \\
\texttt{║◌│◌◌│◌◌║ 7} \\
\texttt{╚═╧══╧══╝} \\
\texttt{```} \\
\texttt{Зараз ви граєте раунд 5 з максимум 12.} \\
\texttt{Разом ви накопичили 9 з 8 штрафів.} \\
\texttt{Інший гравець перемістив об'єкт на своїй сітці.} \\
\\ 
\\ 
\texttt{Яка ваша наступна команда?} \\
            }
        }
    }
    & \\ \\

    \theutterance \stepcounter{utterance}  
    & & & \multicolumn{4}{p{0.6\linewidth}}{
        \cellcolor[rgb]{0.9,0.9,0.9}{
            \makecell[{{p{\linewidth}}}]{
                \texttt{\tiny{[GM$\langle$P2]}}
                \texttt{ПЕРЕСУНЬ: P, (3,5)} \\
            }
        }
    }
     \\ \\

    \theutterance \stepcounter{utterance}  
    & & & \multicolumn{2}{p{0.3\linewidth}}{
        \cellcolor[rgb]{0.9,0.9,0.9}{
            \makecell[{{p{\linewidth}}}]{
                \texttt{\tiny{[GM$|$GM]}}
                \texttt{ПЕРЕСУНЬ: P, (3,5)} \\
            }
        }
    }
    & & \\ \\

    \theutterance \stepcounter{utterance}  
    & & & \multicolumn{2}{p{0.3\linewidth}}{
        \cellcolor[rgb]{0.9,0.9,0.9}{
            \makecell[{{p{\linewidth}}}]{
                \texttt{\tiny{[GM$|$GM]}}
                \texttt{Штраф: (3,5) не порожня, а містить '─'.} \\
\\ 
\texttt{Разом ви накопичили 10 з 8 штрафів.} \\
\texttt{Будь ласка, спробуйте ще раз!} \\
            }
        }
    }
    & & \\ \\

    \theutterance \stepcounter{utterance}  
    & & \multicolumn{4}{p{0.6\linewidth}}{
        \cellcolor[rgb]{0.9,0.9,0.9}{
            \makecell[{{p{\linewidth}}}]{
                \texttt{\tiny{[GM$\rangle$P2]}}
                \texttt{Штраф: (3,5) не порожня, а містить '─'.} \\
\\ 
\texttt{Разом ви накопичили 10 з 8 штрафів.} \\
\texttt{Будь ласка, спробуйте ще раз!} \\
            }
        }
    }
    & \\ \\

    \theutterance \stepcounter{utterance}  
    & & & \multicolumn{4}{p{0.6\linewidth}}{
        \cellcolor[rgb]{0.9,0.9,0.9}{
            \makecell[{{p{\linewidth}}}]{
                \texttt{\tiny{[GM$\langle$P2]}}
                \texttt{ПЕРЕСУНЬ: P, (4,5)} \\
            }
        }
    }
     \\ \\

    \theutterance \stepcounter{utterance}  
    & & & \multicolumn{2}{p{0.3\linewidth}}{
        \cellcolor[rgb]{0.9,0.9,0.9}{
            \makecell[{{p{\linewidth}}}]{
                \texttt{\tiny{[GM$|$GM]}}
                \texttt{ПЕРЕСУНЬ: P, (4,5)} \\
            }
        }
    }
    & & \\ \\

    \theutterance \stepcounter{utterance}  
    & & & \multicolumn{2}{p{0.3\linewidth}}{
        \cellcolor[rgb]{0.9,0.9,0.9}{
            \makecell[{{p{\linewidth}}}]{
                \texttt{\tiny{[GM$|$GM]}}
                \texttt{Штраф: (4,5) не порожня, а містить '┘'.} \\
\\ 
\texttt{Разом ви накопичили 11 з 8 штрафів.} \\
\texttt{Будь ласка, спробуйте ще раз!} \\
            }
        }
    }
    & & \\ \\

    \theutterance \stepcounter{utterance}  
    & & \multicolumn{4}{p{0.6\linewidth}}{
        \cellcolor[rgb]{0.9,0.9,0.9}{
            \makecell[{{p{\linewidth}}}]{
                \texttt{\tiny{[GM$\rangle$P2]}}
                \texttt{Штраф: (4,5) не порожня, а містить '┘'.} \\
\\ 
\texttt{Разом ви накопичили 11 з 8 штрафів.} \\
\texttt{Будь ласка, спробуйте ще раз!} \\
            }
        }
    }
    & \\ \\

    \theutterance \stepcounter{utterance}  
    & & & \multicolumn{4}{p{0.6\linewidth}}{
        \cellcolor[rgb]{0.9,0.9,0.9}{
            \makecell[{{p{\linewidth}}}]{
                \texttt{\tiny{[GM$\langle$P2]}}
                \texttt{ПЕРЕСУНЬ: P, (5,5)} \\
            }
        }
    }
     \\ \\

    \theutterance \stepcounter{utterance}  
    & & & \multicolumn{2}{p{0.3\linewidth}}{
        \cellcolor[rgb]{0.9,0.9,0.9}{
            \makecell[{{p{\linewidth}}}]{
                \texttt{\tiny{[GM$|$GM]}}
                \texttt{ПЕРЕСУНЬ: P, (5,5)} \\
            }
        }
    }
    & & \\ \\

    \theutterance \stepcounter{utterance}  
    & & & \multicolumn{2}{p{0.3\linewidth}}{
        \cellcolor[rgb]{0.9,0.9,0.9}{
            \makecell[{{p{\linewidth}}}]{
                \texttt{\tiny{[GM$|$GM]}}
                \texttt{P moved to (5, 5)} \\
            }
        }
    }
    & & \\ \\

    \theutterance \stepcounter{utterance}  
    & & \multicolumn{4}{p{0.6\linewidth}}{
        \cellcolor[rgb]{0.9,0.9,0.9}{
            \makecell[{{p{\linewidth}}}]{
                \texttt{\tiny{[P1$\langle$GM]}}
                \texttt{Об'єкт 'P' успішно переміщено в (5,5). Ваша оновлена сітка виглядає так:} \\
\texttt{```} \\
\texttt{ 1234567} \\
\texttt{╔═╤═╤═══╗} \\
\texttt{║◌│◌│◌◌◌║ 1} \\
\texttt{║◌├─┤◌◌◌║ 2} \\
\texttt{╟─┤L├─┬─╢ 3} \\
\texttt{║◌│C│◌│◌║ 4} \\
\texttt{╟─┼─┘P│◌║ 5} \\
\texttt{║◌│◌◌┌┴─╢ 6} \\
\texttt{║◌│◌◌│◌◌║ 7} \\
\texttt{╚═╧══╧══╝} \\
\texttt{```} \\
\texttt{Зараз ви граєте раунд 5 з максимум 12.} \\
\texttt{Разом ви накопичили 11 з 8 штрафів.} \\
\texttt{Інший гравець перемістив об'єкт на своїй сітці.} \\
\\ 
\\ 
\texttt{Яка ваша наступна команда?} \\
            }
        }
    }
    & \\ \\

    \theutterance \stepcounter{utterance}  
    & \multicolumn{4}{p{0.6\linewidth}}{
        \cellcolor[rgb]{0.9,0.9,0.9}{
            \makecell[{{p{\linewidth}}}]{
                \texttt{\tiny{[P1$\rangle$GM]}}
                \texttt{СКАЖИ: Я вважаю, що наш поточний прогрес демонструє складність досягнення цієї конфігурації, особливо після штрафів. Можливо, варто переглянути загальну стратегію. Чи потрібно нам уточнити цільові координати об'єктів, щоб уникнути подальших помилок?} \\
            }
        }
    }
    & & \\ \\

    \theutterance \stepcounter{utterance}  
    & & & \multicolumn{2}{p{0.3\linewidth}}{
        \cellcolor[rgb]{0.9,0.9,0.9}{
            \makecell[{{p{\linewidth}}}]{
                \texttt{\tiny{[GM$|$GM]}}
                \texttt{СКАЖИ: Я вважаю, що наш поточний прогрес демонструє складність досягнення цієї конфігурації, особливо після штрафів. Можливо, варто переглянути загальну стратегію. Чи потрібно нам уточнити цільові координати об'єктів, щоб уникнути подальших помилок?} \\
            }
        }
    }
    & & \\ \\

    \theutterance \stepcounter{utterance}  
    & & \multicolumn{4}{p{0.6\linewidth}}{
        \cellcolor[rgb]{0.9,0.9,0.9}{
            \makecell[{{p{\linewidth}}}]{
                \texttt{\tiny{[GM$\rangle$P2]}}
                \texttt{Об'єкт 'P' успішно переміщено в (5,5). Ваша оновлена сітка виглядає так:} \\
\texttt{```} \\
\texttt{ 1234567} \\
\texttt{╔═╤═╤═══╗} \\
\texttt{║◌│◌│◌◌◌║ 1} \\
\texttt{║L├─┤◌◌◌║ 2} \\
\texttt{╟─┤◌├─┬─╢ 3} \\
\texttt{║◌│C│◌│◌║ 4} \\
\texttt{╟─┼─┘P│◌║ 5} \\
\texttt{║◌│◌◌┌┴─╢ 6} \\
\texttt{║◌│◌◌│◌◌║ 7} \\
\texttt{╚═╧══╧══╝} \\
\texttt{```} \\
\texttt{Зараз ви граєте раунд 6 з максимум 12.} \\
\texttt{Разом ви накопичили 11 з 8 штрафів.} \\
\texttt{Інший гравець надіслав це повідомлення:} \\
\\ 
\\ 
\texttt{"Я вважаю, що наш поточний прогрес демонструє складність досягнення цієї конфігурації, особливо після штрафів. Можливо, варто переглянути загальну стратегію. Чи потрібно нам уточнити цільові координати об'єктів, щоб уникнути подальших помилок?"} \\
\\ 
\\ 
\texttt{Яка ваша наступна команда?} \\
            }
        }
    }
    & \\ \\

    \theutterance \stepcounter{utterance}  
    & & & \multicolumn{4}{p{0.6\linewidth}}{
        \cellcolor[rgb]{0.9,0.9,0.9}{
            \makecell[{{p{\linewidth}}}]{
                \texttt{\tiny{[GM$\langle$P2]}}
                \texttt{СКАЖИ: Згоден, страждання показують, що координати треба уточняти. Давайте перевіримо цільові точки ще раз, перш ніж продовжувати.} \\
            }
        }
    }
     \\ \\

    \theutterance \stepcounter{utterance}  
    & & & \multicolumn{2}{p{0.3\linewidth}}{
        \cellcolor[rgb]{0.9,0.9,0.9}{
            \makecell[{{p{\linewidth}}}]{
                \texttt{\tiny{[GM$|$GM]}}
                \texttt{СКАЖИ: Згоден, страждання показують, що координати треба уточняти. Давайте перевіримо цільові точки ще раз, перш ніж продовжувати.} \\
            }
        }
    }
    & & \\ \\

    \theutterance \stepcounter{utterance}  
    & & \multicolumn{4}{p{0.6\linewidth}}{
        \cellcolor[rgb]{0.9,0.9,0.9}{
            \makecell[{{p{\linewidth}}}]{
                \texttt{\tiny{[P1$\langle$GM]}}
                \texttt{Ваше повідомлення передано іншому гравцю.} \\
\\ 
\\ 
\texttt{Зараз ви граєте раунд 6 з максимум 12.} \\
\texttt{Разом ви накопичили 11 з 8 штрафів.} \\
\texttt{Інший гравець надіслав це повідомлення:} \\
\\ 
\\ 
\texttt{"Згоден, страждання показують, що координати треба уточняти. Давайте перевіримо цільові точки ще раз, перш ніж продовжувати."} \\
\\ 
\\ 
\texttt{Яка ваша наступна команда?} \\
            }
        }
    }
    & \\ \\

    \theutterance \stepcounter{utterance}  
    & \multicolumn{4}{p{0.6\linewidth}}{
        \cellcolor[rgb]{0.9,0.9,0.9}{
            \makecell[{{p{\linewidth}}}]{
                \texttt{\tiny{[P1$\rangle$GM]}}
                \texttt{СКАЖИ: Я пропоную перевірити точні поточні позиції кожного об'єкта на обох полях, щоб уточнити цільові координати. Можливо, я помилився в початкових пропозиціях. Який твій варіант узгодження цільових координат?} \\
            }
        }
    }
    & & \\ \\

    \theutterance \stepcounter{utterance}  
    & & & \multicolumn{2}{p{0.3\linewidth}}{
        \cellcolor[rgb]{0.9,0.9,0.9}{
            \makecell[{{p{\linewidth}}}]{
                \texttt{\tiny{[GM$|$GM]}}
                \texttt{СКАЖИ: Я пропоную перевірити точні поточні позиції кожного об'єкта на обох полях, щоб уточнити цільові координати. Можливо, я помилився в початкових пропозиціях. Який твій варіант узгодження цільових координат?} \\
            }
        }
    }
    & & \\ \\

    \theutterance \stepcounter{utterance}  
    & & \multicolumn{4}{p{0.6\linewidth}}{
        \cellcolor[rgb]{0.9,0.9,0.9}{
            \makecell[{{p{\linewidth}}}]{
                \texttt{\tiny{[GM$\rangle$P2]}}
                \texttt{Ваше повідомлення передано іншому гравцю.} \\
\\ 
\\ 
\texttt{Зараз ви граєте раунд 7 з максимум 12.} \\
\texttt{Разом ви накопичили 11 з 8 штрафів.} \\
\texttt{Інший гравець надіслав це повідомлення:} \\
\\ 
\\ 
\texttt{"Я пропоную перевірити точні поточні позиції кожного об'єкта на обох полях, щоб уточнити цільові координати. Можливо, я помилився в початкових пропозиціях. Який твій варіант узгодження цільових координат?"} \\
\\ 
\\ 
\texttt{Яка ваша наступна команда?} \\
            }
        }
    }
    & \\ \\

    \theutterance \stepcounter{utterance}  
    & & & \multicolumn{4}{p{0.6\linewidth}}{
        \cellcolor[rgb]{0.9,0.9,0.9}{
            \makecell[{{p{\linewidth}}}]{
                \texttt{\tiny{[GM$\langle$P2]}}
                \texttt{СКАЖИ: Згоден, перевіримо поточні позиції. Давайте перемістимо всі об'єкти в такі точки: L — (1,3), C — (2,4), P — (3,5). Це зробить наші конфігурації однаковими та простими для вирівнювання.} \\
            }
        }
    }
     \\ \\

    \theutterance \stepcounter{utterance}  
    & & & \multicolumn{2}{p{0.3\linewidth}}{
        \cellcolor[rgb]{0.9,0.9,0.9}{
            \makecell[{{p{\linewidth}}}]{
                \texttt{\tiny{[GM$|$GM]}}
                \texttt{СКАЖИ: Згоден, перевіримо поточні позиції. Давайте перемістимо всі об'єкти в такі точки: L — (1,3), C — (2,4), P — (3,5). Це зробить наші конфігурації однаковими та простими для вирівнювання.} \\
            }
        }
    }
    & & \\ \\

    \theutterance \stepcounter{utterance}  
    & & \multicolumn{4}{p{0.6\linewidth}}{
        \cellcolor[rgb]{0.9,0.9,0.9}{
            \makecell[{{p{\linewidth}}}]{
                \texttt{\tiny{[P1$\langle$GM]}}
                \texttt{Ваше повідомлення передано іншому гравцю.} \\
\\ 
\\ 
\texttt{Зараз ви граєте раунд 7 з максимум 12.} \\
\texttt{Разом ви накопичили 11 з 8 штрафів.} \\
\texttt{Інший гравець надіслав це повідомлення:} \\
\\ 
\\ 
\texttt{"Згоден, перевіримо поточні позиції. Давайте перемістимо всі об'єкти в такі точки: L — (1,3), C — (2,4), P — (3,5). Це зробить наші конфігурації однаковими та простими для вирівнювання."} \\
\\ 
\\ 
\texttt{Яка ваша наступна команда?} \\
            }
        }
    }
    & \\ \\

    \theutterance \stepcounter{utterance}  
    & \multicolumn{4}{p{0.6\linewidth}}{
        \cellcolor[rgb]{0.9,0.9,0.9}{
            \makecell[{{p{\linewidth}}}]{
                \texttt{\tiny{[P1$\rangle$GM]}}
                \texttt{ПЕРЕСУНЬ: L, (1,3)} \\
            }
        }
    }
    & & \\ \\

    \theutterance \stepcounter{utterance}  
    & & & \multicolumn{2}{p{0.3\linewidth}}{
        \cellcolor[rgb]{0.9,0.9,0.9}{
            \makecell[{{p{\linewidth}}}]{
                \texttt{\tiny{[GM$|$GM]}}
                \texttt{ПЕРЕСУНЬ: L, (1,3)} \\
            }
        }
    }
    & & \\ \\

    \theutterance \stepcounter{utterance}  
    & & & \multicolumn{2}{p{0.3\linewidth}}{
        \cellcolor[rgb]{0.9,0.9,0.9}{
            \makecell[{{p{\linewidth}}}]{
                \texttt{\tiny{[GM$|$GM]}}
                \texttt{Штраф: (1,3) не порожня, а містить '─'.} \\
\\ 
\texttt{Разом ви накопичили 12 з 8 штрафів.} \\
\texttt{Будь ласка, спробуйте ще раз!} \\
            }
        }
    }
    & & \\ \\

    \theutterance \stepcounter{utterance}  
    & & \multicolumn{4}{p{0.6\linewidth}}{
        \cellcolor[rgb]{0.9,0.9,0.9}{
            \makecell[{{p{\linewidth}}}]{
                \texttt{\tiny{[P1$\langle$GM]}}
                \texttt{Штраф: (1,3) не порожня, а містить '─'.} \\
\\ 
\texttt{Разом ви накопичили 12 з 8 штрафів.} \\
\texttt{Будь ласка, спробуйте ще раз!} \\
            }
        }
    }
    & \\ \\

    \theutterance \stepcounter{utterance}  
    & \multicolumn{4}{p{0.6\linewidth}}{
        \cellcolor[rgb]{0.9,0.9,0.9}{
            \makecell[{{p{\linewidth}}}]{
                \texttt{\tiny{[P1$\rangle$GM]}}
                \texttt{ПЕРЕСУНЬ: L, (2,3)} \\
            }
        }
    }
    & & \\ \\

    \theutterance \stepcounter{utterance}  
    & & & \multicolumn{2}{p{0.3\linewidth}}{
        \cellcolor[rgb]{0.9,0.9,0.9}{
            \makecell[{{p{\linewidth}}}]{
                \texttt{\tiny{[GM$|$GM]}}
                \texttt{ПЕРЕСУНЬ: L, (2,3)} \\
            }
        }
    }
    & & \\ \\

    \theutterance \stepcounter{utterance}  
    & & & \multicolumn{2}{p{0.3\linewidth}}{
        \cellcolor[rgb]{0.9,0.9,0.9}{
            \makecell[{{p{\linewidth}}}]{
                \texttt{\tiny{[GM$|$GM]}}
                \texttt{Штраф: (2,3) не порожня, а містить '┤'.} \\
\\ 
\texttt{Разом ви накопичили 13 з 8 штрафів.} \\
\texttt{Будь ласка, спробуйте ще раз!} \\
            }
        }
    }
    & & \\ \\

    \theutterance \stepcounter{utterance}  
    & & \multicolumn{4}{p{0.6\linewidth}}{
        \cellcolor[rgb]{0.9,0.9,0.9}{
            \makecell[{{p{\linewidth}}}]{
                \texttt{\tiny{[P1$\langle$GM]}}
                \texttt{Штраф: (2,3) не порожня, а містить '┤'.} \\
\\ 
\texttt{Разом ви накопичили 13 з 8 штрафів.} \\
\texttt{Будь ласка, спробуйте ще раз!} \\
            }
        }
    }
    & \\ \\

    \theutterance \stepcounter{utterance}  
    & \multicolumn{4}{p{0.6\linewidth}}{
        \cellcolor[rgb]{0.9,0.9,0.9}{
            \makecell[{{p{\linewidth}}}]{
                \texttt{\tiny{[P1$\rangle$GM]}}
                \texttt{ПЕРЕСУНЬ: L, (3,3)} \\
            }
        }
    }
    & & \\ \\

    \theutterance \stepcounter{utterance}  
    & & & \multicolumn{2}{p{0.3\linewidth}}{
        \cellcolor[rgb]{0.9,0.9,0.9}{
            \makecell[{{p{\linewidth}}}]{
                \texttt{\tiny{[GM$|$GM]}}
                \texttt{ПЕРЕСУНЬ: L, (3,3)} \\
            }
        }
    }
    & & \\ \\

    \theutterance \stepcounter{utterance}  
    & & & \multicolumn{2}{p{0.3\linewidth}}{
        \cellcolor[rgb]{0.9,0.9,0.9}{
            \makecell[{{p{\linewidth}}}]{
                \texttt{\tiny{[GM$|$GM]}}
                \texttt{Штраф: (3,3) не порожня, а містить 'L'.} \\
\\ 
\texttt{Разом ви накопичили 14 з 8 штрафів.} \\
\texttt{Будь ласка, спробуйте ще раз!} \\
            }
        }
    }
    & & \\ \\

    \theutterance \stepcounter{utterance}  
    & & \multicolumn{4}{p{0.6\linewidth}}{
        \cellcolor[rgb]{0.9,0.9,0.9}{
            \makecell[{{p{\linewidth}}}]{
                \texttt{\tiny{[P1$\langle$GM]}}
                \texttt{Штраф: (3,3) не порожня, а містить 'L'.} \\
\\ 
\texttt{Разом ви накопичили 14 з 8 штрафів.} \\
\texttt{Будь ласка, спробуйте ще раз!} \\
            }
        }
    }
    & \\ \\

    \theutterance \stepcounter{utterance}  
    & \multicolumn{4}{p{0.6\linewidth}}{
        \cellcolor[rgb]{0.9,0.9,0.9}{
            \makecell[{{p{\linewidth}}}]{
                \texttt{\tiny{[P1$\rangle$GM]}}
                \texttt{ПЕРЕСУНЬ: L, (4,3)} \\
            }
        }
    }
    & & \\ \\

    \theutterance \stepcounter{utterance}  
    & & & \multicolumn{2}{p{0.3\linewidth}}{
        \cellcolor[rgb]{0.9,0.9,0.9}{
            \makecell[{{p{\linewidth}}}]{
                \texttt{\tiny{[GM$|$GM]}}
                \texttt{ПЕРЕСУНЬ: L, (4,3)} \\
            }
        }
    }
    & & \\ \\

    \theutterance \stepcounter{utterance}  
    & & & \multicolumn{2}{p{0.3\linewidth}}{
        \cellcolor[rgb]{0.9,0.9,0.9}{
            \makecell[{{p{\linewidth}}}]{
                \texttt{\tiny{[GM$|$GM]}}
                \texttt{Штраф: (4,3) не порожня, а містить '├'.} \\
\\ 
\texttt{Разом ви накопичили 15 з 8 штрафів.} \\
\texttt{Будь ласка, спробуйте ще раз!} \\
            }
        }
    }
    & & \\ \\

    \theutterance \stepcounter{utterance}  
    & & \multicolumn{4}{p{0.6\linewidth}}{
        \cellcolor[rgb]{0.9,0.9,0.9}{
            \makecell[{{p{\linewidth}}}]{
                \texttt{\tiny{[P1$\langle$GM]}}
                \texttt{Штраф: (4,3) не порожня, а містить '├'.} \\
\\ 
\texttt{Разом ви накопичили 15 з 8 штрафів.} \\
\texttt{Будь ласка, спробуйте ще раз!} \\
            }
        }
    }
    & \\ \\

    \theutterance \stepcounter{utterance}  
    & \multicolumn{4}{p{0.6\linewidth}}{
        \cellcolor[rgb]{0.9,0.9,0.9}{
            \makecell[{{p{\linewidth}}}]{
                \texttt{\tiny{[P1$\rangle$GM]}}
                \texttt{ПЕРЕСУНЬ: L, (5,3)} \\
            }
        }
    }
    & & \\ \\

    \theutterance \stepcounter{utterance}  
    & & & \multicolumn{2}{p{0.3\linewidth}}{
        \cellcolor[rgb]{0.9,0.9,0.9}{
            \makecell[{{p{\linewidth}}}]{
                \texttt{\tiny{[GM$|$GM]}}
                \texttt{ПЕРЕСУНЬ: L, (5,3)} \\
            }
        }
    }
    & & \\ \\

    \theutterance \stepcounter{utterance}  
    & & & \multicolumn{2}{p{0.3\linewidth}}{
        \cellcolor[rgb]{0.9,0.9,0.9}{
            \makecell[{{p{\linewidth}}}]{
                \texttt{\tiny{[GM$|$GM]}}
                \texttt{Штраф: (5,3) не порожня, а містить '─'.} \\
\\ 
\texttt{Разом ви накопичили 16 з 8 штрафів.} \\
\texttt{Будь ласка, спробуйте ще раз!} \\
            }
        }
    }
    & & \\ \\

    \theutterance \stepcounter{utterance}  
    & & \multicolumn{4}{p{0.6\linewidth}}{
        \cellcolor[rgb]{0.9,0.9,0.9}{
            \makecell[{{p{\linewidth}}}]{
                \texttt{\tiny{[P1$\langle$GM]}}
                \texttt{Штраф: (5,3) не порожня, а містить '─'.} \\
\\ 
\texttt{Разом ви накопичили 16 з 8 штрафів.} \\
\texttt{Будь ласка, спробуйте ще раз!} \\
            }
        }
    }
    & \\ \\

    \theutterance \stepcounter{utterance}  
    & \multicolumn{4}{p{0.6\linewidth}}{
        \cellcolor[rgb]{0.9,0.9,0.9}{
            \makecell[{{p{\linewidth}}}]{
                \texttt{\tiny{[P1$\rangle$GM]}}
                \texttt{ПЕРЕСУНЬ: L, (6,3)} \\
            }
        }
    }
    & & \\ \\

    \theutterance \stepcounter{utterance}  
    & & & \multicolumn{2}{p{0.3\linewidth}}{
        \cellcolor[rgb]{0.9,0.9,0.9}{
            \makecell[{{p{\linewidth}}}]{
                \texttt{\tiny{[GM$|$GM]}}
                \texttt{ПЕРЕСУНЬ: L, (6,3)} \\
            }
        }
    }
    & & \\ \\

    \theutterance \stepcounter{utterance}  
    & & & \multicolumn{2}{p{0.3\linewidth}}{
        \cellcolor[rgb]{0.9,0.9,0.9}{
            \makecell[{{p{\linewidth}}}]{
                \texttt{\tiny{[GM$|$GM]}}
                \texttt{Штраф: (6,3) не порожня, а містить '┬'.} \\
\\ 
\texttt{Разом ви накопичили 17 з 8 штрафів.} \\
\texttt{Будь ласка, спробуйте ще раз!} \\
            }
        }
    }
    & & \\ \\

    \theutterance \stepcounter{utterance}  
    & & \multicolumn{4}{p{0.6\linewidth}}{
        \cellcolor[rgb]{0.9,0.9,0.9}{
            \makecell[{{p{\linewidth}}}]{
                \texttt{\tiny{[P1$\langle$GM]}}
                \texttt{Штраф: (6,3) не порожня, а містить '┬'.} \\
\\ 
\texttt{Разом ви накопичили 17 з 8 штрафів.} \\
\texttt{Будь ласка, спробуйте ще раз!} \\
            }
        }
    }
    & \\ \\

    \theutterance \stepcounter{utterance}  
    & \multicolumn{4}{p{0.6\linewidth}}{
        \cellcolor[rgb]{0.9,0.9,0.9}{
            \makecell[{{p{\linewidth}}}]{
                \texttt{\tiny{[P1$\rangle$GM]}}
                \texttt{ПЕРЕСУНЬ: L, (7,3)} \\
            }
        }
    }
    & & \\ \\

    \theutterance \stepcounter{utterance}  
    & & & \multicolumn{2}{p{0.3\linewidth}}{
        \cellcolor[rgb]{0.9,0.9,0.9}{
            \makecell[{{p{\linewidth}}}]{
                \texttt{\tiny{[GM$|$GM]}}
                \texttt{ПЕРЕСУНЬ: L, (7,3)} \\
            }
        }
    }
    & & \\ \\

    \theutterance \stepcounter{utterance}  
    & & & \multicolumn{2}{p{0.3\linewidth}}{
        \cellcolor[rgb]{0.9,0.9,0.9}{
            \makecell[{{p{\linewidth}}}]{
                \texttt{\tiny{[GM$|$GM]}}
                \texttt{Штраф: (7,3) не порожня, а містить '─'.} \\
\\ 
\texttt{Разом ви накопичили 18 з 8 штрафів.} \\
\texttt{Будь ласка, спробуйте ще раз!} \\
            }
        }
    }
    & & \\ \\

    \theutterance \stepcounter{utterance}  
    & & \multicolumn{4}{p{0.6\linewidth}}{
        \cellcolor[rgb]{0.9,0.9,0.9}{
            \makecell[{{p{\linewidth}}}]{
                \texttt{\tiny{[P1$\langle$GM]}}
                \texttt{Штраф: (7,3) не порожня, а містить '─'.} \\
\\ 
\texttt{Разом ви накопичили 18 з 8 штрафів.} \\
\texttt{Будь ласка, спробуйте ще раз!} \\
            }
        }
    }
    & \\ \\

    \theutterance \stepcounter{utterance}  
    & \multicolumn{4}{p{0.6\linewidth}}{
        \cellcolor[rgb]{0.9,0.9,0.9}{
            \makecell[{{p{\linewidth}}}]{
                \texttt{\tiny{[P1$\rangle$GM]}}
                \texttt{ПЕРЕСУНЬ: L, (1,4)} \\
            }
        }
    }
    & & \\ \\

    \theutterance \stepcounter{utterance}  
    & & & \multicolumn{2}{p{0.3\linewidth}}{
        \cellcolor[rgb]{0.9,0.9,0.9}{
            \makecell[{{p{\linewidth}}}]{
                \texttt{\tiny{[GM$|$GM]}}
                \texttt{ПЕРЕСУНЬ: L, (1,4)} \\
            }
        }
    }
    & & \\ \\

    \theutterance \stepcounter{utterance}  
    & & & \multicolumn{2}{p{0.3\linewidth}}{
        \cellcolor[rgb]{0.9,0.9,0.9}{
            \makecell[{{p{\linewidth}}}]{
                \texttt{\tiny{[GM$|$GM]}}
                \texttt{L moved to (1, 4)} \\
            }
        }
    }
    & & \\ \\

    \theutterance \stepcounter{utterance}  
    & & \multicolumn{4}{p{0.6\linewidth}}{
        \cellcolor[rgb]{0.9,0.9,0.9}{
            \makecell[{{p{\linewidth}}}]{
                \texttt{\tiny{[GM$\rangle$P2]}}
                \texttt{Ваше повідомлення передано іншому гравцю.} \\
\\ 
\\ 
\texttt{Зараз ви граєте раунд 8 з максимум 12.} \\
\texttt{Разом ви накопичили 18 з 8 штрафів.} \\
\texttt{Інший гравець перемістив об'єкт на своїй сітці.} \\
\\ 
\\ 
\texttt{Яка ваша наступна команда?} \\
            }
        }
    }
    & \\ \\

    \theutterance \stepcounter{utterance}  
    & & & \multicolumn{4}{p{0.6\linewidth}}{
        \cellcolor[rgb]{0.9,0.9,0.9}{
            \makecell[{{p{\linewidth}}}]{
                \texttt{\tiny{[GM$\langle$P2]}}
                \texttt{ПЕРЕСУНЬ: L, (1,3)} \\
            }
        }
    }
     \\ \\

    \theutterance \stepcounter{utterance}  
    & & & \multicolumn{2}{p{0.3\linewidth}}{
        \cellcolor[rgb]{0.9,0.9,0.9}{
            \makecell[{{p{\linewidth}}}]{
                \texttt{\tiny{[GM$|$GM]}}
                \texttt{ПЕРЕСУНЬ: L, (1,3)} \\
            }
        }
    }
    & & \\ \\

    \theutterance \stepcounter{utterance}  
    & & & \multicolumn{2}{p{0.3\linewidth}}{
        \cellcolor[rgb]{0.9,0.9,0.9}{
            \makecell[{{p{\linewidth}}}]{
                \texttt{\tiny{[GM$|$GM]}}
                \texttt{Штраф: (1,3) не порожня, а містить '─'.} \\
\\ 
\texttt{Разом ви накопичили 19 з 8 штрафів.} \\
\texttt{Будь ласка, спробуйте ще раз!} \\
            }
        }
    }
    & & \\ \\

    \theutterance \stepcounter{utterance}  
    & & \multicolumn{4}{p{0.6\linewidth}}{
        \cellcolor[rgb]{0.9,0.9,0.9}{
            \makecell[{{p{\linewidth}}}]{
                \texttt{\tiny{[GM$\rangle$P2]}}
                \texttt{Штраф: (1,3) не порожня, а містить '─'.} \\
\\ 
\texttt{Разом ви накопичили 19 з 8 штрафів.} \\
\texttt{Будь ласка, спробуйте ще раз!} \\
            }
        }
    }
    & \\ \\

    \theutterance \stepcounter{utterance}  
    & & & \multicolumn{4}{p{0.6\linewidth}}{
        \cellcolor[rgb]{0.9,0.9,0.9}{
            \makecell[{{p{\linewidth}}}]{
                \texttt{\tiny{[GM$\langle$P2]}}
                \texttt{ПЕРЕСУНЬ: L, (1,4)} \\
            }
        }
    }
     \\ \\

    \theutterance \stepcounter{utterance}  
    & & & \multicolumn{2}{p{0.3\linewidth}}{
        \cellcolor[rgb]{0.9,0.9,0.9}{
            \makecell[{{p{\linewidth}}}]{
                \texttt{\tiny{[GM$|$GM]}}
                \texttt{ПЕРЕСУНЬ: L, (1,4)} \\
            }
        }
    }
    & & \\ \\

    \theutterance \stepcounter{utterance}  
    & & & \multicolumn{2}{p{0.3\linewidth}}{
        \cellcolor[rgb]{0.9,0.9,0.9}{
            \makecell[{{p{\linewidth}}}]{
                \texttt{\tiny{[GM$|$GM]}}
                \texttt{L moved to (1, 4)} \\
            }
        }
    }
    & & \\ \\

    \theutterance \stepcounter{utterance}  
    & & \multicolumn{4}{p{0.6\linewidth}}{
        \cellcolor[rgb]{0.9,0.9,0.9}{
            \makecell[{{p{\linewidth}}}]{
                \texttt{\tiny{[P1$\langle$GM]}}
                \texttt{Об'єкт 'L' успішно переміщено в (1,4). Ваша оновлена сітка виглядає так:} \\
\texttt{```} \\
\texttt{ 1234567} \\
\texttt{╔═╤═╤═══╗} \\
\texttt{║◌│◌│◌◌◌║ 1} \\
\texttt{║◌├─┤◌◌◌║ 2} \\
\texttt{╟─┤◌├─┬─╢ 3} \\
\texttt{║L│C│◌│◌║ 4} \\
\texttt{╟─┼─┘P│◌║ 5} \\
\texttt{║◌│◌◌┌┴─╢ 6} \\
\texttt{║◌│◌◌│◌◌║ 7} \\
\texttt{╚═╧══╧══╝} \\
\texttt{```} \\
\texttt{Зараз ви граєте раунд 8 з максимум 12.} \\
\texttt{Разом ви накопичили 19 з 8 штрафів.} \\
\texttt{Інший гравець перемістив об'єкт на своїй сітці.} \\
\\ 
\\ 
\texttt{Яка ваша наступна команда?} \\
            }
        }
    }
    & \\ \\

    \theutterance \stepcounter{utterance}  
    & \multicolumn{4}{p{0.6\linewidth}}{
        \cellcolor[rgb]{0.9,0.9,0.9}{
            \makecell[{{p{\linewidth}}}]{
                \texttt{\tiny{[P1$\rangle$GM]}}
                \texttt{СКАЖИ: Я помітив, що наші цілі координати все ще не враховують успішне переміщення L у (1,4). Подайте, будь ласка, оновлені цільові координати для L, C та P, пов'язані з оновленою вашою сіткою.} \\
            }
        }
    }
    & & \\ \\

    \theutterance \stepcounter{utterance}  
    & & & \multicolumn{2}{p{0.3\linewidth}}{
        \cellcolor[rgb]{0.9,0.9,0.9}{
            \makecell[{{p{\linewidth}}}]{
                \texttt{\tiny{[GM$|$GM]}}
                \texttt{СКАЖИ: Я помітив, що наші цілі координати все ще не враховують успішне переміщення L у (1,4). Подайте, будь ласка, оновлені цільові координати для L, C та P, пов'язані з оновленою вашою сіткою.} \\
            }
        }
    }
    & & \\ \\

    \theutterance \stepcounter{utterance}  
    & & \multicolumn{4}{p{0.6\linewidth}}{
        \cellcolor[rgb]{0.9,0.9,0.9}{
            \makecell[{{p{\linewidth}}}]{
                \texttt{\tiny{[GM$\rangle$P2]}}
                \texttt{Об'єкт 'L' успішно переміщено в (1,4). Ваша оновлена сітка виглядає так:} \\
\texttt{```} \\
\texttt{ 1234567} \\
\texttt{╔═╤═╤═══╗} \\
\texttt{║◌│◌│◌◌◌║ 1} \\
\texttt{║◌├─┤◌◌◌║ 2} \\
\texttt{╟─┤◌├─┬─╢ 3} \\
\texttt{║L│C│◌│◌║ 4} \\
\texttt{╟─┼─┘P│◌║ 5} \\
\texttt{║◌│◌◌┌┴─╢ 6} \\
\texttt{║◌│◌◌│◌◌║ 7} \\
\texttt{╚═╧══╧══╝} \\
\texttt{```} \\
\texttt{Зараз ви граєте раунд 9 з максимум 12.} \\
\texttt{Разом ви накопичили 19 з 8 штрафів.} \\
\texttt{Інший гравець надіслав це повідомлення:} \\
\\ 
\\ 
\texttt{"Я помітив, що наші цілі координати все ще не враховують успішне переміщення L у (1,4). Подайте, будь ласка, оновлені цільові координати для L, C та P, пов'язані з оновленою вашою сіткою."} \\
\\ 
\\ 
\texttt{Яка ваша наступна команда?} \\
            }
        }
    }
    & \\ \\

    \theutterance \stepcounter{utterance}  
    & & & \multicolumn{4}{p{0.6\linewidth}}{
        \cellcolor[rgb]{0.9,0.9,0.9}{
            \makecell[{{p{\linewidth}}}]{
                \texttt{\tiny{[GM$\langle$P2]}}
                \texttt{СКАЖИ: Згідно з оновленою сіткою, нові цілі координати: L — (1,4), C — (2,5), P — (3,6). Це точна відповідність.} \\
            }
        }
    }
     \\ \\

    \theutterance \stepcounter{utterance}  
    & & & \multicolumn{2}{p{0.3\linewidth}}{
        \cellcolor[rgb]{0.9,0.9,0.9}{
            \makecell[{{p{\linewidth}}}]{
                \texttt{\tiny{[GM$|$GM]}}
                \texttt{СКАЖИ: Згідно з оновленою сіткою, нові цілі координати: L — (1,4), C — (2,5), P — (3,6). Це точна відповідність.} \\
            }
        }
    }
    & & \\ \\

    \theutterance \stepcounter{utterance}  
    & & \multicolumn{4}{p{0.6\linewidth}}{
        \cellcolor[rgb]{0.9,0.9,0.9}{
            \makecell[{{p{\linewidth}}}]{
                \texttt{\tiny{[P1$\langle$GM]}}
                \texttt{Ваше повідомлення передано іншому гравцю.} \\
\\ 
\\ 
\texttt{Зараз ви граєте раунд 9 з максимум 12.} \\
\texttt{Разом ви накопичили 19 з 8 штрафів.} \\
\texttt{Інший гравець надіслав це повідомлення:} \\
\\ 
\\ 
\texttt{"Згідно з оновленою сіткою, нові цілі координати: L — (1,4), C — (2,5), P — (3,6). Це точна відповідність."} \\
\\ 
\\ 
\texttt{Яка ваша наступна команда?} \\
            }
        }
    }
    & \\ \\

    \theutterance \stepcounter{utterance}  
    & \multicolumn{4}{p{0.6\linewidth}}{
        \cellcolor[rgb]{0.9,0.9,0.9}{
            \makecell[{{p{\linewidth}}}]{
                \texttt{\tiny{[P1$\rangle$GM]}}
                \texttt{ПЕРЕСУНЬ: C, (2,5)} \\
            }
        }
    }
    & & \\ \\

    \theutterance \stepcounter{utterance}  
    & & & \multicolumn{2}{p{0.3\linewidth}}{
        \cellcolor[rgb]{0.9,0.9,0.9}{
            \makecell[{{p{\linewidth}}}]{
                \texttt{\tiny{[GM$|$GM]}}
                \texttt{ПЕРЕСУНЬ: C, (2,5)} \\
            }
        }
    }
    & & \\ \\

    \theutterance \stepcounter{utterance}  
    & & & \multicolumn{2}{p{0.3\linewidth}}{
        \cellcolor[rgb]{0.9,0.9,0.9}{
            \makecell[{{p{\linewidth}}}]{
                \texttt{\tiny{[GM$|$GM]}}
                \texttt{Штраф: (2,5) не порожня, а містить '┼'.} \\
\\ 
\texttt{Разом ви накопичили 20 з 8 штрафів.} \\
\texttt{Будь ласка, спробуйте ще раз!} \\
            }
        }
    }
    & & \\ \\

    \theutterance \stepcounter{utterance}  
    & & \multicolumn{4}{p{0.6\linewidth}}{
        \cellcolor[rgb]{0.9,0.9,0.9}{
            \makecell[{{p{\linewidth}}}]{
                \texttt{\tiny{[P1$\langle$GM]}}
                \texttt{Штраф: (2,5) не порожня, а містить '┼'.} \\
\\ 
\texttt{Разом ви накопичили 20 з 8 штрафів.} \\
\texttt{Будь ласка, спробуйте ще раз!} \\
            }
        }
    }
    & \\ \\

    \theutterance \stepcounter{utterance}  
    & \multicolumn{4}{p{0.6\linewidth}}{
        \cellcolor[rgb]{0.9,0.9,0.9}{
            \makecell[{{p{\linewidth}}}]{
                \texttt{\tiny{[P1$\rangle$GM]}}
                \texttt{ПЕРЕСУНЬ: C, (3,5)} \\
            }
        }
    }
    & & \\ \\

    \theutterance \stepcounter{utterance}  
    & & & \multicolumn{2}{p{0.3\linewidth}}{
        \cellcolor[rgb]{0.9,0.9,0.9}{
            \makecell[{{p{\linewidth}}}]{
                \texttt{\tiny{[GM$|$GM]}}
                \texttt{ПЕРЕСУНЬ: C, (3,5)} \\
            }
        }
    }
    & & \\ \\

    \theutterance \stepcounter{utterance}  
    & & & \multicolumn{2}{p{0.3\linewidth}}{
        \cellcolor[rgb]{0.9,0.9,0.9}{
            \makecell[{{p{\linewidth}}}]{
                \texttt{\tiny{[GM$|$GM]}}
                \texttt{Штраф: (3,5) не порожня, а містить '─'.} \\
\\ 
\texttt{Разом ви накопичили 21 з 8 штрафів.} \\
\texttt{Будь ласка, спробуйте ще раз!} \\
            }
        }
    }
    & & \\ \\

    \theutterance \stepcounter{utterance}  
    & & \multicolumn{4}{p{0.6\linewidth}}{
        \cellcolor[rgb]{0.9,0.9,0.9}{
            \makecell[{{p{\linewidth}}}]{
                \texttt{\tiny{[P1$\langle$GM]}}
                \texttt{Штраф: (3,5) не порожня, а містить '─'.} \\
\\ 
\texttt{Разом ви накопичили 21 з 8 штрафів.} \\
\texttt{Будь ласка, спробуйте ще раз!} \\
            }
        }
    }
    & \\ \\

    \theutterance \stepcounter{utterance}  
    & \multicolumn{4}{p{0.6\linewidth}}{
        \cellcolor[rgb]{0.9,0.9,0.9}{
            \makecell[{{p{\linewidth}}}]{
                \texttt{\tiny{[P1$\rangle$GM]}}
                \texttt{ПЕРЕСУНЬ: C, (4,5)} \\
            }
        }
    }
    & & \\ \\

    \theutterance \stepcounter{utterance}  
    & & & \multicolumn{2}{p{0.3\linewidth}}{
        \cellcolor[rgb]{0.9,0.9,0.9}{
            \makecell[{{p{\linewidth}}}]{
                \texttt{\tiny{[GM$|$GM]}}
                \texttt{ПЕРЕСУНЬ: C, (4,5)} \\
            }
        }
    }
    & & \\ \\

    \theutterance \stepcounter{utterance}  
    & & & \multicolumn{2}{p{0.3\linewidth}}{
        \cellcolor[rgb]{0.9,0.9,0.9}{
            \makecell[{{p{\linewidth}}}]{
                \texttt{\tiny{[GM$|$GM]}}
                \texttt{Штраф: (4,5) не порожня, а містить '┘'.} \\
\\ 
\texttt{Разом ви накопичили 22 з 8 штрафів.} \\
\texttt{Будь ласка, спробуйте ще раз!} \\
            }
        }
    }
    & & \\ \\

    \theutterance \stepcounter{utterance}  
    & & \multicolumn{4}{p{0.6\linewidth}}{
        \cellcolor[rgb]{0.9,0.9,0.9}{
            \makecell[{{p{\linewidth}}}]{
                \texttt{\tiny{[P1$\langle$GM]}}
                \texttt{Штраф: (4,5) не порожня, а містить '┘'.} \\
\\ 
\texttt{Разом ви накопичили 22 з 8 штрафів.} \\
\texttt{Будь ласка, спробуйте ще раз!} \\
            }
        }
    }
    & \\ \\

    \theutterance \stepcounter{utterance}  
    & \multicolumn{4}{p{0.6\linewidth}}{
        \cellcolor[rgb]{0.9,0.9,0.9}{
            \makecell[{{p{\linewidth}}}]{
                \texttt{\tiny{[P1$\rangle$GM]}}
                \texttt{ПЕРЕСУНЬ: C, (5,5)} \\
            }
        }
    }
    & & \\ \\

    \theutterance \stepcounter{utterance}  
    & & & \multicolumn{2}{p{0.3\linewidth}}{
        \cellcolor[rgb]{0.9,0.9,0.9}{
            \makecell[{{p{\linewidth}}}]{
                \texttt{\tiny{[GM$|$GM]}}
                \texttt{ПЕРЕСУНЬ: C, (5,5)} \\
            }
        }
    }
    & & \\ \\

    \theutterance \stepcounter{utterance}  
    & & & \multicolumn{2}{p{0.3\linewidth}}{
        \cellcolor[rgb]{0.9,0.9,0.9}{
            \makecell[{{p{\linewidth}}}]{
                \texttt{\tiny{[GM$|$GM]}}
                \texttt{Штраф: (5,5) не порожня, а містить 'P'.} \\
\\ 
\texttt{Разом ви накопичили 23 з 8 штрафів.} \\
\texttt{Будь ласка, спробуйте ще раз!} \\
            }
        }
    }
    & & \\ \\

    \theutterance \stepcounter{utterance}  
    & & \multicolumn{4}{p{0.6\linewidth}}{
        \cellcolor[rgb]{0.9,0.9,0.9}{
            \makecell[{{p{\linewidth}}}]{
                \texttt{\tiny{[P1$\langle$GM]}}
                \texttt{Штраф: (5,5) не порожня, а містить 'P'.} \\
\\ 
\texttt{Разом ви накопичили 23 з 8 штрафів.} \\
\texttt{Будь ласка, спробуйте ще раз!} \\
            }
        }
    }
    & \\ \\

    \theutterance \stepcounter{utterance}  
    & \multicolumn{4}{p{0.6\linewidth}}{
        \cellcolor[rgb]{0.9,0.9,0.9}{
            \makecell[{{p{\linewidth}}}]{
                \texttt{\tiny{[P1$\rangle$GM]}}
                \texttt{ПЕРЕСУНЬ: C, (6,5)} \\
            }
        }
    }
    & & \\ \\

    \theutterance \stepcounter{utterance}  
    & & & \multicolumn{2}{p{0.3\linewidth}}{
        \cellcolor[rgb]{0.9,0.9,0.9}{
            \makecell[{{p{\linewidth}}}]{
                \texttt{\tiny{[GM$|$GM]}}
                \texttt{ПЕРЕСУНЬ: C, (6,5)} \\
            }
        }
    }
    & & \\ \\

    \theutterance \stepcounter{utterance}  
    & & & \multicolumn{2}{p{0.3\linewidth}}{
        \cellcolor[rgb]{0.9,0.9,0.9}{
            \makecell[{{p{\linewidth}}}]{
                \texttt{\tiny{[GM$|$GM]}}
                \texttt{Штраф: (6,5) не порожня, а містить '│'.} \\
\\ 
\texttt{Разом ви накопичили 24 з 8 штрафів.} \\
\texttt{Будь ласка, спробуйте ще раз!} \\
            }
        }
    }
    & & \\ \\

    \theutterance \stepcounter{utterance}  
    & & \multicolumn{4}{p{0.6\linewidth}}{
        \cellcolor[rgb]{0.9,0.9,0.9}{
            \makecell[{{p{\linewidth}}}]{
                \texttt{\tiny{[P1$\langle$GM]}}
                \texttt{Штраф: (6,5) не порожня, а містить '│'.} \\
\\ 
\texttt{Разом ви накопичили 24 з 8 штрафів.} \\
\texttt{Будь ласка, спробуйте ще раз!} \\
            }
        }
    }
    & \\ \\

    \theutterance \stepcounter{utterance}  
    & \multicolumn{4}{p{0.6\linewidth}}{
        \cellcolor[rgb]{0.9,0.9,0.9}{
            \makecell[{{p{\linewidth}}}]{
                \texttt{\tiny{[P1$\rangle$GM]}}
                \texttt{ПЕРЕСУНЬ: C, (7,5)} \\
            }
        }
    }
    & & \\ \\

    \theutterance \stepcounter{utterance}  
    & & & \multicolumn{2}{p{0.3\linewidth}}{
        \cellcolor[rgb]{0.9,0.9,0.9}{
            \makecell[{{p{\linewidth}}}]{
                \texttt{\tiny{[GM$|$GM]}}
                \texttt{ПЕРЕСУНЬ: C, (7,5)} \\
            }
        }
    }
    & & \\ \\

    \theutterance \stepcounter{utterance}  
    & & & \multicolumn{2}{p{0.3\linewidth}}{
        \cellcolor[rgb]{0.9,0.9,0.9}{
            \makecell[{{p{\linewidth}}}]{
                \texttt{\tiny{[GM$|$GM]}}
                \texttt{C moved to (7, 5)} \\
            }
        }
    }
    & & \\ \\

    \theutterance \stepcounter{utterance}  
    & & \multicolumn{4}{p{0.6\linewidth}}{
        \cellcolor[rgb]{0.9,0.9,0.9}{
            \makecell[{{p{\linewidth}}}]{
                \texttt{\tiny{[GM$\rangle$P2]}}
                \texttt{Ваше повідомлення передано іншому гравцю.} \\
\\ 
\\ 
\texttt{Зараз ви граєте раунд 10 з максимум 12.} \\
\texttt{Разом ви накопичили 24 з 8 штрафів.} \\
\texttt{Інший гравець перемістив об'єкт на своїй сітці.} \\
\\ 
\\ 
\texttt{Яка ваша наступна команда?} \\
            }
        }
    }
    & \\ \\

    \theutterance \stepcounter{utterance}  
    & & & \multicolumn{4}{p{0.6\linewidth}}{
        \cellcolor[rgb]{0.9,0.9,0.9}{
            \makecell[{{p{\linewidth}}}]{
                \texttt{\tiny{[GM$\langle$P2]}}
                \texttt{ПЕРЕСУНЬ: C, (2,5)} \\
            }
        }
    }
     \\ \\

    \theutterance \stepcounter{utterance}  
    & & & \multicolumn{2}{p{0.3\linewidth}}{
        \cellcolor[rgb]{0.9,0.9,0.9}{
            \makecell[{{p{\linewidth}}}]{
                \texttt{\tiny{[GM$|$GM]}}
                \texttt{ПЕРЕСУНЬ: C, (2,5)} \\
            }
        }
    }
    & & \\ \\

    \theutterance \stepcounter{utterance}  
    & & & \multicolumn{2}{p{0.3\linewidth}}{
        \cellcolor[rgb]{0.9,0.9,0.9}{
            \makecell[{{p{\linewidth}}}]{
                \texttt{\tiny{[GM$|$GM]}}
                \texttt{Штраф: (2,5) не порожня, а містить '┼'.} \\
\\ 
\texttt{Разом ви накопичили 25 з 8 штрафів.} \\
\texttt{Будь ласка, спробуйте ще раз!} \\
            }
        }
    }
    & & \\ \\

    \theutterance \stepcounter{utterance}  
    & & \multicolumn{4}{p{0.6\linewidth}}{
        \cellcolor[rgb]{0.9,0.9,0.9}{
            \makecell[{{p{\linewidth}}}]{
                \texttt{\tiny{[GM$\rangle$P2]}}
                \texttt{Штраф: (2,5) не порожня, а містить '┼'.} \\
\\ 
\texttt{Разом ви накопичили 25 з 8 штрафів.} \\
\texttt{Будь ласка, спробуйте ще раз!} \\
            }
        }
    }
    & \\ \\

    \theutterance \stepcounter{utterance}  
    & & & \multicolumn{4}{p{0.6\linewidth}}{
        \cellcolor[rgb]{0.9,0.9,0.9}{
            \makecell[{{p{\linewidth}}}]{
                \texttt{\tiny{[GM$\langle$P2]}}
                \texttt{ПЕРЕСУНЬ: C, (3,5)} \\
            }
        }
    }
     \\ \\

    \theutterance \stepcounter{utterance}  
    & & & \multicolumn{2}{p{0.3\linewidth}}{
        \cellcolor[rgb]{0.9,0.9,0.9}{
            \makecell[{{p{\linewidth}}}]{
                \texttt{\tiny{[GM$|$GM]}}
                \texttt{ПЕРЕСУНЬ: C, (3,5)} \\
            }
        }
    }
    & & \\ \\

    \theutterance \stepcounter{utterance}  
    & & & \multicolumn{2}{p{0.3\linewidth}}{
        \cellcolor[rgb]{0.9,0.9,0.9}{
            \makecell[{{p{\linewidth}}}]{
                \texttt{\tiny{[GM$|$GM]}}
                \texttt{Штраф: (3,5) не порожня, а містить '─'.} \\
\\ 
\texttt{Разом ви накопичили 26 з 8 штрафів.} \\
\texttt{Будь ласка, спробуйте ще раз!} \\
            }
        }
    }
    & & \\ \\

    \theutterance \stepcounter{utterance}  
    & & \multicolumn{4}{p{0.6\linewidth}}{
        \cellcolor[rgb]{0.9,0.9,0.9}{
            \makecell[{{p{\linewidth}}}]{
                \texttt{\tiny{[GM$\rangle$P2]}}
                \texttt{Штраф: (3,5) не порожня, а містить '─'.} \\
\\ 
\texttt{Разом ви накопичили 26 з 8 штрафів.} \\
\texttt{Будь ласка, спробуйте ще раз!} \\
            }
        }
    }
    & \\ \\

    \theutterance \stepcounter{utterance}  
    & & & \multicolumn{4}{p{0.6\linewidth}}{
        \cellcolor[rgb]{0.9,0.9,0.9}{
            \makecell[{{p{\linewidth}}}]{
                \texttt{\tiny{[GM$\langle$P2]}}
                \texttt{ПЕРЕСУНЬ: C, (1,5)} \\
            }
        }
    }
     \\ \\

    \theutterance \stepcounter{utterance}  
    & & & \multicolumn{2}{p{0.3\linewidth}}{
        \cellcolor[rgb]{0.9,0.9,0.9}{
            \makecell[{{p{\linewidth}}}]{
                \texttt{\tiny{[GM$|$GM]}}
                \texttt{ПЕРЕСУНЬ: C, (1,5)} \\
            }
        }
    }
    & & \\ \\

    \theutterance \stepcounter{utterance}  
    & & & \multicolumn{2}{p{0.3\linewidth}}{
        \cellcolor[rgb]{0.9,0.9,0.9}{
            \makecell[{{p{\linewidth}}}]{
                \texttt{\tiny{[GM$|$GM]}}
                \texttt{Штраф: (1,5) не порожня, а містить '─'.} \\
\\ 
\texttt{Разом ви накопичили 27 з 8 штрафів.} \\
\texttt{Будь ласка, спробуйте ще раз!} \\
            }
        }
    }
    & & \\ \\

    \theutterance \stepcounter{utterance}  
    & & \multicolumn{4}{p{0.6\linewidth}}{
        \cellcolor[rgb]{0.9,0.9,0.9}{
            \makecell[{{p{\linewidth}}}]{
                \texttt{\tiny{[GM$\rangle$P2]}}
                \texttt{Штраф: (1,5) не порожня, а містить '─'.} \\
\\ 
\texttt{Разом ви накопичили 27 з 8 штрафів.} \\
\texttt{Будь ласка, спробуйте ще раз!} \\
            }
        }
    }
    & \\ \\

    \theutterance \stepcounter{utterance}  
    & & & \multicolumn{4}{p{0.6\linewidth}}{
        \cellcolor[rgb]{0.9,0.9,0.9}{
            \makecell[{{p{\linewidth}}}]{
                \texttt{\tiny{[GM$\langle$P2]}}
                \texttt{ПЕРЕСУНЬ: C, (4,5)} \\
            }
        }
    }
     \\ \\

    \theutterance \stepcounter{utterance}  
    & & & \multicolumn{2}{p{0.3\linewidth}}{
        \cellcolor[rgb]{0.9,0.9,0.9}{
            \makecell[{{p{\linewidth}}}]{
                \texttt{\tiny{[GM$|$GM]}}
                \texttt{ПЕРЕСУНЬ: C, (4,5)} \\
            }
        }
    }
    & & \\ \\

    \theutterance \stepcounter{utterance}  
    & & & \multicolumn{2}{p{0.3\linewidth}}{
        \cellcolor[rgb]{0.9,0.9,0.9}{
            \makecell[{{p{\linewidth}}}]{
                \texttt{\tiny{[GM$|$GM]}}
                \texttt{Штраф: (4,5) не порожня, а містить '┘'.} \\
\\ 
\texttt{Разом ви накопичили 28 з 8 штрафів.} \\
\texttt{Будь ласка, спробуйте ще раз!} \\
            }
        }
    }
    & & \\ \\

    \theutterance \stepcounter{utterance}  
    & & \multicolumn{4}{p{0.6\linewidth}}{
        \cellcolor[rgb]{0.9,0.9,0.9}{
            \makecell[{{p{\linewidth}}}]{
                \texttt{\tiny{[GM$\rangle$P2]}}
                \texttt{Штраф: (4,5) не порожня, а містить '┘'.} \\
\\ 
\texttt{Разом ви накопичили 28 з 8 штрафів.} \\
\texttt{Будь ласка, спробуйте ще раз!} \\
            }
        }
    }
    & \\ \\

    \theutterance \stepcounter{utterance}  
    & & & \multicolumn{4}{p{0.6\linewidth}}{
        \cellcolor[rgb]{0.9,0.9,0.9}{
            \makecell[{{p{\linewidth}}}]{
                \texttt{\tiny{[GM$\langle$P2]}}
                \texttt{ПЕРЕСУНЬ: C, (2,4)} \\
            }
        }
    }
     \\ \\

    \theutterance \stepcounter{utterance}  
    & & & \multicolumn{2}{p{0.3\linewidth}}{
        \cellcolor[rgb]{0.9,0.9,0.9}{
            \makecell[{{p{\linewidth}}}]{
                \texttt{\tiny{[GM$|$GM]}}
                \texttt{ПЕРЕСУНЬ: C, (2,4)} \\
            }
        }
    }
    & & \\ \\

    \theutterance \stepcounter{utterance}  
    & & & \multicolumn{2}{p{0.3\linewidth}}{
        \cellcolor[rgb]{0.9,0.9,0.9}{
            \makecell[{{p{\linewidth}}}]{
                \texttt{\tiny{[GM$|$GM]}}
                \texttt{Штраф: (2,4) не порожня, а містить '│'.} \\
\\ 
\texttt{Разом ви накопичили 29 з 8 штрафів.} \\
\texttt{Будь ласка, спробуйте ще раз!} \\
            }
        }
    }
    & & \\ \\

    \theutterance \stepcounter{utterance}  
    & & \multicolumn{4}{p{0.6\linewidth}}{
        \cellcolor[rgb]{0.9,0.9,0.9}{
            \makecell[{{p{\linewidth}}}]{
                \texttt{\tiny{[GM$\rangle$P2]}}
                \texttt{Штраф: (2,4) не порожня, а містить '│'.} \\
\\ 
\texttt{Разом ви накопичили 29 з 8 штрафів.} \\
\texttt{Будь ласка, спробуйте ще раз!} \\
            }
        }
    }
    & \\ \\

    \theutterance \stepcounter{utterance}  
    & & & \multicolumn{4}{p{0.6\linewidth}}{
        \cellcolor[rgb]{0.9,0.9,0.9}{
            \makecell[{{p{\linewidth}}}]{
                \texttt{\tiny{[GM$\langle$P2]}}
                \texttt{ПЕРЕСУНЬ: C, (3,4)} \\
            }
        }
    }
     \\ \\

    \theutterance \stepcounter{utterance}  
    & & & \multicolumn{2}{p{0.3\linewidth}}{
        \cellcolor[rgb]{0.9,0.9,0.9}{
            \makecell[{{p{\linewidth}}}]{
                \texttt{\tiny{[GM$|$GM]}}
                \texttt{ПЕРЕСУНЬ: C, (3,4)} \\
            }
        }
    }
    & & \\ \\

    \theutterance \stepcounter{utterance}  
    & & & \multicolumn{2}{p{0.3\linewidth}}{
        \cellcolor[rgb]{0.9,0.9,0.9}{
            \makecell[{{p{\linewidth}}}]{
                \texttt{\tiny{[GM$|$GM]}}
                \texttt{Штраф: (3,4) не порожня, а містить 'C'.} \\
\\ 
\texttt{Разом ви накопичили 30 з 8 штрафів.} \\
\texttt{Будь ласка, спробуйте ще раз!} \\
            }
        }
    }
    & & \\ \\

    \theutterance \stepcounter{utterance}  
    & & \multicolumn{4}{p{0.6\linewidth}}{
        \cellcolor[rgb]{0.9,0.9,0.9}{
            \makecell[{{p{\linewidth}}}]{
                \texttt{\tiny{[GM$\rangle$P2]}}
                \texttt{Штраф: (3,4) не порожня, а містить 'C'.} \\
\\ 
\texttt{Разом ви накопичили 30 з 8 штрафів.} \\
\texttt{Будь ласка, спробуйте ще раз!} \\
            }
        }
    }
    & \\ \\

    \theutterance \stepcounter{utterance}  
    & & & \multicolumn{4}{p{0.6\linewidth}}{
        \cellcolor[rgb]{0.9,0.9,0.9}{
            \makecell[{{p{\linewidth}}}]{
                \texttt{\tiny{[GM$\langle$P2]}}
                \texttt{ПЕРЕСУНЬ: C, (5,4)} \\
            }
        }
    }
     \\ \\

    \theutterance \stepcounter{utterance}  
    & & & \multicolumn{2}{p{0.3\linewidth}}{
        \cellcolor[rgb]{0.9,0.9,0.9}{
            \makecell[{{p{\linewidth}}}]{
                \texttt{\tiny{[GM$|$GM]}}
                \texttt{ПЕРЕСУНЬ: C, (5,4)} \\
            }
        }
    }
    & & \\ \\

    \theutterance \stepcounter{utterance}  
    & & & \multicolumn{2}{p{0.3\linewidth}}{
        \cellcolor[rgb]{0.9,0.9,0.9}{
            \makecell[{{p{\linewidth}}}]{
                \texttt{\tiny{[GM$|$GM]}}
                \texttt{C moved to (5, 4)} \\
            }
        }
    }
    & & \\ \\

    \theutterance \stepcounter{utterance}  
    & & \multicolumn{4}{p{0.6\linewidth}}{
        \cellcolor[rgb]{0.9,0.9,0.9}{
            \makecell[{{p{\linewidth}}}]{
                \texttt{\tiny{[P1$\langle$GM]}}
                \texttt{Об'єкт 'C' успішно переміщено в (7,5). Ваша оновлена сітка виглядає так:} \\
\texttt{```} \\
\texttt{ 1234567} \\
\texttt{╔═╤═╤═══╗} \\
\texttt{║◌│◌│◌◌◌║ 1} \\
\texttt{║◌├─┤◌◌◌║ 2} \\
\texttt{╟─┤◌├─┬─╢ 3} \\
\texttt{║L│◌│◌│◌║ 4} \\
\texttt{╟─┼─┘P│C║ 5} \\
\texttt{║◌│◌◌┌┴─╢ 6} \\
\texttt{║◌│◌◌│◌◌║ 7} \\
\texttt{╚═╧══╧══╝} \\
\texttt{```} \\
\texttt{Зараз ви граєте раунд 10 з максимум 12.} \\
\texttt{Разом ви накопичили 30 з 8 штрафів.} \\
\texttt{Інший гравець перемістив об'єкт на своїй сітці.} \\
\\ 
\\ 
\texttt{Яка ваша наступна команда?} \\
            }
        }
    }
    & \\ \\

    \theutterance \stepcounter{utterance}  
    & \multicolumn{4}{p{0.6\linewidth}}{
        \cellcolor[rgb]{0.9,0.9,0.9}{
            \makecell[{{p{\linewidth}}}]{
                \texttt{\tiny{[P1$\rangle$GM]}}
                \texttt{СКАЖИ: Я вважаю, що наш поточний статус з 30 штрафами показує, що потрібно почати з нуля та планувати кожне переміщення ретельно. Давайте спочатку перевіримо стан нашого та іншого поля щодо поточної конфігурації L, C та P, щоб узгодити нову стратегію мінімізації помилок.} \\
            }
        }
    }
    & & \\ \\

    \theutterance \stepcounter{utterance}  
    & & & \multicolumn{2}{p{0.3\linewidth}}{
        \cellcolor[rgb]{0.9,0.9,0.9}{
            \makecell[{{p{\linewidth}}}]{
                \texttt{\tiny{[GM$|$GM]}}
                \texttt{СКАЖИ: Я вважаю, що наш поточний статус з 30 штрафами показує, що потрібно почати з нуля та планувати кожне переміщення ретельно. Давайте спочатку перевіримо стан нашого та іншого поля щодо поточної конфігурації L, C та P, щоб узгодити нову стратегію мінімізації помилок.} \\
            }
        }
    }
    & & \\ \\

    \theutterance \stepcounter{utterance}  
    & & \multicolumn{4}{p{0.6\linewidth}}{
        \cellcolor[rgb]{0.9,0.9,0.9}{
            \makecell[{{p{\linewidth}}}]{
                \texttt{\tiny{[GM$\rangle$P2]}}
                \texttt{Об'єкт 'C' успішно переміщено в (5,4). Ваша оновлена сітка виглядає так:} \\
\texttt{```} \\
\texttt{ 1234567} \\
\texttt{╔═╤═╤═══╗} \\
\texttt{║◌│◌│◌◌◌║ 1} \\
\texttt{║◌├─┤◌◌◌║ 2} \\
\texttt{╟─┤◌├─┬─╢ 3} \\
\texttt{║L│◌│C│◌║ 4} \\
\texttt{╟─┼─┘P│◌║ 5} \\
\texttt{║◌│◌◌┌┴─╢ 6} \\
\texttt{║◌│◌◌│◌◌║ 7} \\
\texttt{╚═╧══╧══╝} \\
\texttt{```} \\
\texttt{Зараз ви граєте раунд 11 з максимум 12.} \\
\texttt{Разом ви накопичили 30 з 8 штрафів.} \\
\texttt{Інший гравець надіслав це повідомлення:} \\
\\ 
\\ 
\texttt{"Я вважаю, що наш поточний статус з 30 штрафами показує, що потрібно почати з нуля та планувати кожне переміщення ретельно. Давайте спочатку перевіримо стан нашого та іншого поля щодо поточної конфігурації L, C та P, щоб узгодити нову стратегію мінімізації помилок."} \\
\\ 
\\ 
\texttt{Яка ваша наступна команда?} \\
            }
        }
    }
    & \\ \\

    \theutterance \stepcounter{utterance}  
    & & & \multicolumn{4}{p{0.6\linewidth}}{
        \cellcolor[rgb]{0.9,0.9,0.9}{
            \makecell[{{p{\linewidth}}}]{
                \texttt{\tiny{[GM$\langle$P2]}}
                \texttt{СКАЖИ: Згоден, ми маємо значні штрафи і потрібно переоцінити стратегію з нуля. Давайте опишемо поточну позицію нашого поля іншому гравцю та узгодимо детальний план, що включатиме лише перевірені цільові координати для кожного об'єкта одночасно.} \\
            }
        }
    }
     \\ \\

    \theutterance \stepcounter{utterance}  
    & & & \multicolumn{2}{p{0.3\linewidth}}{
        \cellcolor[rgb]{0.9,0.9,0.9}{
            \makecell[{{p{\linewidth}}}]{
                \texttt{\tiny{[GM$|$GM]}}
                \texttt{СКАЖИ: Згоден, ми маємо значні штрафи і потрібно переоцінити стратегію з нуля. Давайте опишемо поточну позицію нашого поля іншому гравцю та узгодимо детальний план, що включатиме лише перевірені цільові координати для кожного об'єкта одночасно.} \\
            }
        }
    }
    & & \\ \\

    \theutterance \stepcounter{utterance}  
    & & \multicolumn{4}{p{0.6\linewidth}}{
        \cellcolor[rgb]{0.9,0.9,0.9}{
            \makecell[{{p{\linewidth}}}]{
                \texttt{\tiny{[P1$\langle$GM]}}
                \texttt{Ваше повідомлення передано іншому гравцю.} \\
\\ 
\\ 
\texttt{Зараз ви граєте раунд 11 з максимум 12.} \\
\texttt{Разом ви накопичили 30 з 8 штрафів.} \\
\texttt{Інший гравець надіслав це повідомлення:} \\
\\ 
\\ 
\texttt{"Згоден, ми маємо значні штрафи і потрібно переоцінити стратегію з нуля. Давайте опишемо поточну позицію нашого поля іншому гравцю та узгодимо детальний план, що включатиме лише перевірені цільові координати для кожного об'єкта одночасно."} \\
\\ 
\\ 
\texttt{Яка ваша наступна команда?} \\
            }
        }
    }
    & \\ \\

    \theutterance \stepcounter{utterance}  
    & \multicolumn{4}{p{0.6\linewidth}}{
        \cellcolor[rgb]{0.9,0.9,0.9}{
            \makecell[{{p{\linewidth}}}]{
                \texttt{\tiny{[P1$\rangle$GM]}}
                \texttt{СКАЖИ: Я розумію, наша низька ефективність через штрафи вимагає докладного опису. Ось поточна конфігурація моєї сітки: L у (1,4), C у (7,5), P у (5,6). Я пропоную іншому гравцю повідомити про їх поточне розташування на його сітці, щоб ми обидва могли миттєво адаптувати наші стратегії та координати для точного вирівнювання.} \\
            }
        }
    }
    & & \\ \\

    \theutterance \stepcounter{utterance}  
    & & & \multicolumn{2}{p{0.3\linewidth}}{
        \cellcolor[rgb]{0.9,0.9,0.9}{
            \makecell[{{p{\linewidth}}}]{
                \texttt{\tiny{[GM$|$GM]}}
                \texttt{СКАЖИ: Я розумію, наша низька ефективність через штрафи вимагає докладного опису. Ось поточна конфігурація моєї сітки: L у (1,4), C у (7,5), P у (5,6). Я пропоную іншому гравцю повідомити про їх поточне розташування на його сітці, щоб ми обидва могли миттєво адаптувати наші стратегії та координати для точного вирівнювання.} \\
            }
        }
    }
    & & \\ \\

    \theutterance \stepcounter{utterance}  
    & & \multicolumn{4}{p{0.6\linewidth}}{
        \cellcolor[rgb]{0.9,0.9,0.9}{
            \makecell[{{p{\linewidth}}}]{
                \texttt{\tiny{[GM$\rangle$P2]}}
                \texttt{Ваше повідомлення передано іншому гравцю.} \\
\\ 
\\ 
\texttt{Зараз ви граєте раунд 12 з максимум 12.} \\
\texttt{Разом ви накопичили 30 з 8 штрафів.} \\
\texttt{Інший гравець надіслав це повідомлення:} \\
\\ 
\\ 
\texttt{"Я розумію, наша низька ефективність через штрафи вимагає докладного опису. Ось поточна конфігурація моєї сітки: L у (1,4), C у (7,5), P у (5,6). Я пропоную іншому гравцю повідомити про їх поточне розташування на його сітці, щоб ми обидва могли миттєво адаптувати наші стратегії та координати для точного вирівнювання."} \\
\\ 
\\ 
\texttt{Яка ваша наступна команда?} \\
            }
        }
    }
    & \\ \\

    \theutterance \stepcounter{utterance}  
    & & & \multicolumn{4}{p{0.6\linewidth}}{
        \cellcolor[rgb]{0.9,0.9,0.9}{
            \makecell[{{p{\linewidth}}}]{
                \texttt{\tiny{[GM$\langle$P2]}}
                \texttt{СКАЖИ: Повідомляю вам проточну конфігурацію мого поля: L у (1,4), C у (5,4), P у (3,5). Пропоную скористатися цими даними для швидкого налаштування стратегії та координат для обох сіток.} \\
            }
        }
    }
     \\ \\

    \theutterance \stepcounter{utterance}  
    & & & \multicolumn{2}{p{0.3\linewidth}}{
        \cellcolor[rgb]{0.9,0.9,0.9}{
            \makecell[{{p{\linewidth}}}]{
                \texttt{\tiny{[GM$|$GM]}}
                \texttt{СКАЖИ: Повідомляю вам проточну конфігурацію мого поля: L у (1,4), C у (5,4), P у (3,5). Пропоную скористатися цими даними для швидкого налаштування стратегії та координат для обох сіток.} \\
            }
        }
    }
    & & \\ \\

    \theutterance \stepcounter{utterance}  
    & & & \multicolumn{2}{p{0.3\linewidth}}{
        \cellcolor[rgb]{0.9,0.9,0.9}{
            \makecell[{{p{\linewidth}}}]{
                \texttt{\tiny{[GM$|$GM]}}
                \texttt{Maximum number of rounds reached} \\
            }
        }
    }
    & & \\ \\

    \theutterance \stepcounter{utterance}  
    & & & \multicolumn{2}{p{0.3\linewidth}}{
        \cellcolor[rgb]{0.9,0.9,0.9}{
            \makecell[{{p{\linewidth}}}]{
                \texttt{\tiny{[GM$|$GM]}}
                \texttt{Initial states:} \\
\texttt{```} \\
\texttt{Player 1:} \\
\texttt{ 1234567} \\
\texttt{╔═╤═╤═══╗} \\
\texttt{║L│◌│◌◌◌║ 1} \\
\texttt{║C├─┤◌◌◌║ 2} \\
\texttt{╟─┤◌├─┬─╢ 3} \\
\texttt{║◌│◌│◌│◌║ 4} \\
\texttt{╟─┼─┘◌│◌║ 5} \\
\texttt{║◌│◌◌┌┴─╢ 6} \\
\texttt{║◌│◌◌│◌P║ 7} \\
\texttt{╚═╧══╧══╝} \\
\\ 
\\ 
\texttt{Player 2:} \\
\texttt{ 1234567} \\
\texttt{╔═╤═╤═══╗} \\
\texttt{║◌│◌│◌◌◌║ 1} \\
\texttt{║◌├─┤◌L◌║ 2} \\
\texttt{╟─┤P├─┬─╢ 3} \\
\texttt{║◌│◌│◌│◌║ 4} \\
\texttt{╟─┼─┘C│◌║ 5} \\
\texttt{║◌│◌◌┌┴─╢ 6} \\
\texttt{║◌│◌◌│◌◌║ 7} \\
\texttt{╚═╧══╧══╝} \\
\\ 
\texttt{```} \\
            }
        }
    }
    & & \\ \\

    \theutterance \stepcounter{utterance}  
    & & & \multicolumn{2}{p{0.3\linewidth}}{
        \cellcolor[rgb]{0.9,0.9,0.9}{
            \makecell[{{p{\linewidth}}}]{
                \texttt{\tiny{[GM$|$GM]}}
                \texttt{End states:} \\
\texttt{```} \\
\texttt{Player 1:} \\
\texttt{ 1234567} \\
\texttt{╔═╤═╤═══╗} \\
\texttt{║◌│◌│◌◌◌║ 1} \\
\texttt{║◌├─┤◌◌◌║ 2} \\
\texttt{╟─┤◌├─┬─╢ 3} \\
\texttt{║L│◌│◌│◌║ 4} \\
\texttt{╟─┼─┘P│C║ 5} \\
\texttt{║◌│◌◌┌┴─╢ 6} \\
\texttt{║◌│◌◌│◌◌║ 7} \\
\texttt{╚═╧══╧══╝} \\
\\ 
\\ 
\texttt{Player 2:} \\
\texttt{ 1234567} \\
\texttt{╔═╤═╤═══╗} \\
\texttt{║◌│◌│◌◌◌║ 1} \\
\texttt{║◌├─┤◌◌◌║ 2} \\
\texttt{╟─┤◌├─┬─╢ 3} \\
\texttt{║L│◌│C│◌║ 4} \\
\texttt{╟─┼─┘P│◌║ 5} \\
\texttt{║◌│◌◌┌┴─╢ 6} \\
\texttt{║◌│◌◌│◌◌║ 7} \\
\texttt{╚═╧══╧══╝} \\
\\ 
\texttt{```} \\
            }
        }
    }
    & & \\ \\

    \theutterance \stepcounter{utterance}  
    & & & \multicolumn{2}{p{0.3\linewidth}}{
        \cellcolor[rgb]{0.9,0.9,0.9}{
            \makecell[{{p{\linewidth}}}]{
                \texttt{\tiny{[GM$|$GM]}}
                \texttt{* success: True} \\
\texttt{* lose: False} \\
\texttt{* aborted: False} \\
\texttt{{-}{-}{-}{-}{-}{-}{-}} \\
\texttt{* Shifts: 4.00} \\
\texttt{* Max Shifts: 4.00} \\
\texttt{* Min Shifts: 2.00} \\
\texttt{* End Distance Sum: 2.24} \\
\texttt{* Init Distance Sum: 15.76} \\
\texttt{* Expected Distance Sum: 12.57} \\
\texttt{* Penalties: 30.00} \\
\texttt{* Max Penalties: 8.00} \\
\texttt{* Rounds: 12.00} \\
\texttt{* Max Rounds: 12.00} \\
\texttt{* Object Count: 3.00} \\
            }
        }
    }
    & & \\ \\

\end{supertabular}
}

\end{document}
